\documentclass[11pt]{amsart}
\usepackage[T1]{fontenc}
\usepackage[utf8]{inputenc}
\usepackage[british]{babel}
\usepackage[margin=1in]{geometry}
\parindent=8pt
\parskip=3pt
\linespread{1.1}
\usepackage{amsmath,amssymb,amsthm}
\usepackage{palatino,euler,microtype}
% Euler uses the upright shape name n.
\renewcommand{\updefault}{n}
\usepackage{booktabs,longtable,array}
\usepackage{graphicx}
\usepackage{tikz}
\usetikzlibrary{positioning,fit,backgrounds,calc,shapes.geometric}
\input{../figures/jfef-grafos.tikzstyles}
\usepackage{xcolor}
\usepackage{float}
\usepackage[section]{placeins}
\usepackage{needspace}
\definecolor{darkblue}{rgb}{0,0,0.7}
\definecolor{darkred}{rgb}{0.7,0,0}
\newcommand{\defi}[1]{\textit{\color{blue}#1}}
\usepackage[pagebackref=true]{hyperref}
\hypersetup{colorlinks=true,linkcolor=darkred,citecolor=darkblue,urlcolor=darkblue,
pdftitle={Supplement: Linear-size graphs at every strict dismantling level}}
\setlength{\emergencystretch}{2em}


\newcommand{\D}{\mathcal D}
\newcommand{\Cl}{\operatorname{Cl}}
\newcommand{\card}[1]{\lvert#1\rvert}
\newcommand{\nb}[2]{N_{#1}(#2)}
\newcommand{\Nb}[2]{#1[N_{#1}(#2)]}
\newtheorem{theorem}{Theorem}[section]
\newtheorem{lemma}[theorem]{Lemma}
\newtheorem{proposition}[theorem]{Proposition}
\newtheorem{corollary}[theorem]{Corollary}
\newtheorem{definition}[theorem]{Definition}
\newtheorem{remark}[theorem]{Remark}
\renewcommand{\urladdrname}{\textit{ORCID}}
\title[Constructions and finite verifications]{Constructions and finite verifications\\for strict dismantling levels}
\author[J. F. Espinoza]{Jes\'us F. Espinoza\textsuperscript{*}}
\thanks{\textsuperscript{*}Corresponding author}
\author[M. E. Fr\'ias-Armenta]{Mart\'in Eduardo Fr\'ias-Armenta}
\author[R. G. Hern\'andez-Amador]{Rosal\'ia G. Hern\'andez-Amador}
\author[H. A. Hern\'andez-Hern\'andez]{H\'ector A. Hern\'andez-Hern\'andez}
\address{Departamento de Matem\'aticas, Universidad de Sonora, Hermosillo, Sonora, M\'exico}
\email[Jes\'us F. Espinoza]{jesusfrancisco.espinoza@unison.mx}
\urladdr[Jes\'us F. Espinoza]{https://orcid.org/0000-0001-5601-3634}
\email[Mart\'in Eduardo Fr\'ias-Armenta]{eduardo.frias@unison.mx}
\urladdr[Mart\'in Eduardo Fr\'ias-Armenta]{https://orcid.org/0000-0002-7687-1270}
\email[Rosal\'ia G. Hern\'andez-Amador]{rosalia.hernandez@unison.mx}
\urladdr[Rosal\'ia G. Hern\'andez-Amador]{https://orcid.org/0000-0001-9657-7012}
\email[H\'ector A. Hern\'andez-Hern\'andez]{hector.hernandez@unison.mx}
\urladdr[H\'ector A. Hern\'andez-Hern\'andez]{https://orcid.org/0000-0002-4046-4841}
\date{}
\subjclass[2020]{Primary 05C75; Secondary 05C69, 05E45}
\keywords{$k$-dismantlable graph, dismantling hierarchy, clique complex, Euler characteristic, graph construction}

\newcommand{\F}{\mathbb F_2}
\newcommand{\Ht}{\widetilde H}
\newcommand{\Ct}{\widetilde C}
\newcommand{\Lk}[2]{#1[N_{#1}(#2)]}
\newcommand{\lk}{\operatorname{lk}}
\newcommand{\runcheck}[1]{\par\smallskip\noindent\textbf{Reproduce this calculation.}
From the repository directory, run\par\nobreak
\begin{center}\texttt{python reproduce.py --only #1}\end{center}\par\smallskip}
\begin{document}
\begin{abstract}
The recursive dismantling hierarchy measures how a graph can be reduced to one
vertex by testing the graphs induced by successive neighbourhoods. This
supplement supplies complete seed data and finite proofs for families of
strict level $k$ with $4k+4$ and $6k+6$ vertices. We relate constructions with
two auxiliary vertices to the graph obtained by deleting edges between the
exclusive parts of their profiles, and identify the one-edge completion that
produces the eleven-vertex seed of the $3k+8$ family. Finally, we show that
adding a layer shifts the reduced homology of every non-central neighbourhood
by one degree. The accompanying calculations verify the finite hypotheses,
clique counts and homological identities used in these arguments.
\end{abstract}
\maketitle
\raggedbottom
\begin{center}\small
Repository: \url{https://github.com/Geometric-and-Combinatorial-Structures/strictly-k-dismantlable-graphs}\\
Release \texttt{v1.0.0}; file checksums in \texttt{SHA256SUMS.txt}.
\end{center}
\tableofcontents
\section{The constructions and their finite data}
A graph is dismantled by deleting vertices until one vertex remains. In an
ordinary dismantling, each deleted vertex is dominated by a neighbour. The
recursive hierarchy permits a deletion at level $k$ when the graph induced by
its current neighbours is dismantlable at level $k-1$. A graph has
\emph{strict level $k$} if it admits a reduction at level $k$ and none at level
$k-1$.

The article \emph{Linear-size graphs at every strict dismantling level}
constructs graphs $F_k$ and $H_k$ of strict level $k$, for every $k\geq1$,
with
\[
 |V(F_k)|=3k+8,\qquad |V(H_k)|=4k+4.
\]
The title, authors and current article reference are recorded in
\href{https://github.com/Geometric-and-Combinatorial-Structures/strictly-k-dismantlable-graphs/blob/v1.1.1/ARTICLE.md}{\texttt{ARTICLE.md}}.
This supplement develops the finite proof for $H_k$, gives a third family
$J_k$ of order $6k+6$, and explains the kernel completion and homological
identities behind the layer construction. Let $m_k$ be the least order
of a graph of strict level $k$. The two families give
$m_k\leq\min\{4k+4,3k+8\}$; they do not determine $m_k$ in general.
References here to sections, figures and results of the article concern the
identified 23-page manuscript; \texttt{ARTICLE.md} distinguishes it from
the version at the public article locator.

Each family starts with a finite seed. A new layer contains a centre whose
neighbourhood graph is the previous graph, together with auxiliary vertices
that exclude first deletions at the lower level. To apply this construction,
we need a reduction of the seed at its stated level, an obstruction to the
preceding level, a reduction after deleting the first centre, and obstructions
in the remaining neighbourhoods. We supply the finite data for these tasks
and use the homological argument of Section~\ref{sec:homology} to justify the
criterion for every subsequent layer.

For example, $H_1$ has eight vertices. Its reduction at level one requires
only two outer deletions, each justified by a single fold in a neighbourhood;
the remaining graph is a cone. None of its original neighbourhoods is a cone,
so no level-zero reduction can begin. Section~\ref{sec:family4} gives this
example in full before treating the larger recursive certificates.

\noindent\textbf{Organisation.}
Sections~\ref{sec:prelim} and~\ref{sec:construction} fix the hierarchy and
recall the common construction and its Euler criterion. Section~\ref{sec:family4} supplies
the $4k+4$ family, including its two prefixes; Section~\ref{sec:family6}
checks a second application with five auxiliary types. Section~\ref{sec:kernels}
relates constructions with two auxiliary types to an edge-deleted kernel and
identifies the completion producing $F_1$. Its first subsection supplies the
homological preliminaries for Sections~\ref{sec:homology} and~\ref{sec:minimum}.
Section~\ref{sec:homology} proves
the degree shift in non-central neighbourhoods and the resulting homological
criterion. Section~\ref{sec:minimum} explains the finite verification that the minimum at
strict level two is twelve. Appendix~\ref{sec:article} records the finite calculations and
figures for the article's $3k+8$ family.

\noindent\textbf{Conventions and notation.}
The layer indices form $\mathcal L_m$ (written $L_m$ in the article); clique counts are
written $\kappa_t$, reserving $c_i$ for centres. The symbols $F_k$ and $H_k$
retain their meanings in the article; $J_k$ denotes the $6k+6$ family.
Within a section, $A$ denotes its seed, $B$ its first extension and $G_m$ the
graph with $m$ added layers; this layer count is distinct from the minimum
order $m_k$. The seed vertex $a_i$ has numerical label $i$ in
the data files. Homology is
reduced homology over $\mathbb F_2$, with augmented chains in degree $-1$.
The command following each finite calculation runs from the root of the
accompanying repository and requires Python 3.11 or later.

\section{The dismantling hierarchy}\label{sec:prelim}
All graphs are finite, simple and undirected. For $S\subseteq V(H)$, write
$H[S]$ for the induced subgraph and $H-v=H[V(H)\setminus\{v\}]$.
The open and closed neighbourhood sets are
\[
 N_H(v)=\{u:uv\in E(H)\},\qquad N_H[v]=N_H(v)\cup\{v\}.
\]
We write $L_H(v)=H[N_H(v)]$ for the \defi{neighbourhood graph}.
A \defi{cone} is a non-empty graph with a universal vertex, called an
\defi{apex}. Thus $K_1$ is a cone and the empty graph is not.

A vertex $u$ is \defi{dominated} by $w\ne u$ when
$N_H[u]\subseteq N_H[w]$. Deleting it is a \defi{fold}, written $u\to w$.
All neighbourhoods in a sequence of folds are taken in the graph then remaining.

\begin{lemma}\label{lem:dom}
A vertex $u$ is dominated if and only if $L_H(u)$ is a cone.
\end{lemma}
\begin{proof}
The inclusion $N_H[u]\subseteq N_H[w]$ says exactly that $w$ is a
neighbour of $u$ adjacent to all its other neighbours.
\end{proof}

The hierarchy of Fieux and Jouve~\cite{FJ2020} replaces the cone test by
a reduction in the preceding class. At level one these are the
$s$-dismantling moves of~\cite{BFJ2010}.
\begin{definition}\label{def:hierarchy}
Let $\D_{-1}$ consist of the non-empty cones. For $j\geq0$, a graph is
in $\D_j$ if it reduces to $K_1$ by successive deletions whose current
neighbourhood graphs belong to $\D_{j-1}$. Such a graph has
\defi{strict level $j$} when it is outside $\D_{j-1}$.
The empty graph belongs to none of these classes.
\end{definition}
The empty reduction puts $K_1$ in every class. Membership is invariant under
isomorphism, since an isomorphism preserves every neighbourhood test.

\begin{lemma}\label{lem:inclusions}
The classes are nested: $\D_{-1}\subseteq\D_0\subseteq\D_1\subseteq\cdots$.
\end{lemma}
\begin{proof}
A cone folds onto its apex. If $\D_{j-1}\subseteq\D_j$, every
neighbourhood in a level-$j$ reduction also passes the level-$(j+1)$ test.
The same sequence proves the next inclusion.
\end{proof}

\begin{lemma}\label{lem:first}
For $|V(H)|>1$ and $j\geq0$,
\begin{equation}\label{eq:first}
 H\in\D_j\quad\Longleftrightarrow\quad
 \text{some }v\in V(H)\text{ satisfies }L_H(v)\in\D_{j-1}
 \text{ and }H-v\in\D_j.
\end{equation}
\end{lemma}
\begin{proof}
Read the first deletion of a reduction for the forward implication.
For the converse, prepend the permitted deletion of $v$ to a reduction of $H-v$.
\end{proof}
In particular, failure of the neighbourhood test at every vertex excludes
membership. Equation~\eqref{eq:first}, together with the empty graph and cone
cases, is also the finite recursion used by the verification programs.

\section{The common layer construction}\label{sec:construction}
The three families use the same adjacency rules. We recall them here so that
the seed data in the next sections determine each graph without consulting
another document. A centre will recover the preceding graph as its
neighbourhood; auxiliary vertices will provide the obstructions that make
the next level strict.

\subsection{Seed data and adjacency rules}
Choose a non-empty seed $A$, an integer $d\geq1$, and a graph $R$ on the
type set $\mathcal T=\{0,\ldots,d-1\}$. For each type $j$, choose a
\defi{profile} $P_j\subseteq V(A)$: these are its neighbours in the seed.
Let $\mathcal L_m=\{0,\ldots,m-1\}$, with $\mathcal L_0=\varnothing$.
The graph $G_m$ has vertex set
\[
 V(A)\sqcup\{c_i:i\in\mathcal L_m\}
 \sqcup\{x_{j,i}:j\in\mathcal T,\ i\in\mathcal L_m\}.
\]
Here $c_i$ is the centre of layer $i$ and $x_{j,i}$ is its auxiliary of type $j$.
The following rules specify all edges.
\begin{enumerate}
\item The seed induces $A$.
\item The centres form a clique and are adjacent to every seed vertex.
\item All copies of a fixed auxiliary type form a clique.
\item Auxiliaries of distinct types $j,t$ are adjacent exactly when $jt\in E(R)$.
\item A seed vertex $a$ is adjacent to $x_{j,i}$ exactly when $a\in P_j$.
\item The centre $c_i$ is adjacent to $x_{j,h}$ exactly when $i\ne h$.
\end{enumerate}
Put $B=G_1$, $p=c_0$ and $x_j=x_{j,0}$. Then $G_0=A$, $L_B(p)=A$, and
\begin{equation}\label{eq:general-order}
 |V(G_m)|=|V(A)|+(d+1)m.
\end{equation}
A vertex other than a centre is called \defi{non-central}.

The following two facts supply the positive reductions and will also be used
in the homological arguments.
\begin{lemma}\label{lem:centre}
For $m\geq1$ and $i\in\mathcal L_m$, $L_{G_m}(c_i)\cong G_{m-1}$.
\end{lemma}
\begin{proof}
Its vertices are the seed and all layers except layer $i$. Fix the seed and
renumber the remaining layers in increasing order. All six rules are
preserved, including the inequality of layer indices in the last rule.
\end{proof}

\begin{lemma}\label{lem:remainder}
For $m\geq2$, the graph $G_m-c_i$ belongs to $\D_0$.
\end{lemma}
\begin{proof}
For each type $j$ and $h\ne i$, fold $x_{j,h}$ onto $x_{j,i}$. These
vertices have the same neighbours in the seed and among auxiliary types;
after $c_i$ is deleted, the target sees all remaining centres and the source
omits $c_h$. Hence the required closed-neighbourhood inclusion holds.
The inclusions persist after preceding folds, and every target is retained.
After all folds, each remaining centre is universal. At least one exists
because $m\geq2$, so the resulting cone folds to a vertex.
\end{proof}
For $m=1$ there is no remaining centre. The reduction of $B-p$ must be
checked separately for each seed.

\subsection{Clique counts and the Euler criterion}
Let $\kappa_t(H)$ count cliques of cardinality $t$, with $\kappa_0(H)=1$.
The \defi{clique polynomial} and its alternating sum are
\[
 C_H(z)=\sum_{t\geq0}\kappa_t(H)z^t,\qquad q(H)=C_H(-1).
\]
The \defi{clique complex} $\Cl(H)$ consists of the cliques, including the empty
face. With reduced Euler characteristic, including the empty graph convention,
\begin{equation}\label{eq:euler}
 q(H)=1-\chi(\Cl(H))=-\widetilde\chi(\Cl(H)),\qquad q(\varnothing)=1.
\end{equation}
\begin{lemma}\label{lem:qrec}
For every vertex $v$,
\[
 C_H(z)=C_{H-v}(z)+zC_{L_H(v)}(z),\qquad
 q(H)=q(H-v)-q(L_H(v)).
\]
A cone has $q=0$, and a fold preserves $q$.
\end{lemma}
\begin{proof}
Separate cliques according to whether they contain $v$~\cite{HL1994}.
In a cone, adjoining the apex pairs cliques of opposite signs. For a fold,
Lemma~\ref{lem:dom} makes the neighbourhood term vanish.
\end{proof}
Taking coefficients gives
\begin{equation}\label{eq:clique-coefficients}
 \kappa_t(H)=\kappa_t(H-v)+\kappa_{t-1}(L_H(v))\quad(t\geq1).
\end{equation}
The initial values are $\kappa_0(H)=1$ and $\kappa_t(\varnothing)=0$ for $t\geq1$.

\begin{lemma}\label{lem:qobstruction}
Membership in any $\D_j$, $j\geq-1$, implies $q=0$.
\end{lemma}
\begin{proof}
Start with cones. Inductively, all neighbourhood terms in a level-$j$
reduction vanish. The recurrence preserves $q$ down to $K_1$, where it is zero.
\end{proof}
Thus a non-zero $q$ excludes membership; a zero value does not establish it.

\begin{theorem}\label{thm:template}
Suppose that, for some $s\geq0$,
\begin{enumerate}
\item $A\in\D_s\setminus\D_{s-1}$;
\item $B-p\in\D_s$;
\item $q(L_B(v))\ne0$ for every $v\ne p$.
\end{enumerate}
Then $G_m\in\D_{s+m}\setminus\D_{s+m-1}$ for every $m\geq0$, with
order given by~\eqref{eq:general-order}.
\end{theorem}
This is the Euler criterion used in the article. A proof within this supplement
follows from Theorem~\ref{hom:thm:criterion}: by~\eqref{hom:eq:euler}, a non-zero
$q$ gives non-zero reduced homology, and $\D_s\subseteq\D_{s+1}$ supplies
its remainder hypothesis. Section~\ref{sec:homology} proves that theorem
from the adjacency rules, Lemmas~\ref{lem:centre} and~\ref{lem:remainder},
the first-deletion equivalence~\eqref{eq:first}, and the elementary homology
of Section~\ref{sec:homological-tools}.
No other result of Sections~\ref{sec:family4}--\ref{sec:kernels} is used.

\begin{remark}\label{rem:rest-hypothesis}
The same argument allows $B-p\in\D_{s+1}$ in place of condition (2).
It does not remove the requirement to reduce $B-p$.
\end{remark}

The article also proves that these Euler hypotheses require at least two
auxiliary types, hence at least three new vertices per layer. This restriction
concerns the stated template. It is attained by $F_k$ and is not a lower bound
on the order of every graph of strict level $k$.

\section{The family of order \texorpdfstring{$4k+4$}{4k+4}}\label{sec:family4}
The second family in the article begins with graphs on eight, twelve and
sixteen vertices at strict levels one, two and three. All three occur as
nested induced subgraphs. From the sixteen-vertex seed onwards, the layer
construction adds four vertices at a time. We specify these graphs, justify
their finite reductions and obstructions, and then apply the general theorem.

\subsection{The sixteen-vertex seed and its profiles}\label{4k4:sec:seed}
We use three auxiliary types, so each layer adds four vertices. Let
\(V(A)=\{a_0,\ldots,a_{15}\}\), with the adjacency list in
Table~\ref{4k4:tab:Aadj}. It is symmetric, has no loops, and has 97 edges.
The profiles in Table~\ref{4k4:tab:profiles} have sizes 13, 14 and 15.
The type graph is
\begin{equation}\label{4k4:eq:R}
E(R)=\{\{0,1\},\{0,2\}\}.
\end{equation}
\begin{table}[H]
\centering\small
\caption{Complete adjacency list of the sixteen-vertex seed \(A\).}
\label{4k4:tab:Aadj}
\begin{tabular}{r l}\toprule
\(v\)&\(N_A(a_v)\), given by indices\\\midrule
0 & 2, 3, 5, 6, 7, 8, 9, 10, 11, 12, 13, 14\\
1 & 2, 3, 4, 6, 7, 8, 9, 10, 11, 12, 13, 14\\
2 & 0, 1, 3, 6, 7, 8, 9, 10, 11, 12, 13, 15\\
3 & 0, 1, 2, 4, 5, 8, 9, 10, 11, 12, 13, 15\\
4 & 1, 3, 5, 6, 7, 8, 10, 11, 12, 13, 14, 15\\
5 & 0, 3, 4, 6, 7, 9, 10, 11, 12, 13, 14, 15\\
6 & 0, 1, 2, 4, 5, 8, 9, 11, 12, 13, 14, 15\\
7 & 0, 1, 2, 4, 5, 8, 9, 10, 12, 13, 14, 15\\
8 & 0, 1, 2, 3, 4, 6, 7, 9, 10, 12, 14, 15\\
9 & 0, 1, 2, 3, 5, 6, 7, 8, 11, 13, 14, 15\\
10 & 0, 1, 2, 3, 4, 5, 7, 8, 11, 13, 14, 15\\
11 & 0, 1, 2, 3, 4, 5, 6, 9, 10, 12, 14, 15\\
12 & 0, 1, 2, 3, 4, 5, 6, 7, 8, 11, 14, 15\\
13 & 0, 1, 2, 3, 4, 5, 6, 7, 9, 10, 14, 15\\
14 & 0, 1, 4, 5, 6, 7, 8, 9, 10, 11, 12, 13, 15\\
15 & 2, 3, 4, 5, 6, 7, 8, 9, 10, 11, 12, 13, 14\\
\bottomrule\end{tabular}
\end{table}

\begin{table}[H]
\centering
\caption{Profiles specified by their omitted seed vertices. Every copy
of type \(j\) is adjacent to all remaining vertices of \(A\).}
\label{4k4:tab:profiles}
\begin{tabular}{r l r}\toprule
\(j\)&\(V(A)\setminus P_j\)&\(\card{P_j}\)\\\midrule
0 & \(\{a_8,a_{11},a_{14}\}\)&13\\
1 & \(\{a_{11},a_{15}\}\)&14\\
2 & \(\{a_{13}\}\)&15\\
\bottomrule\end{tabular}
\end{table}

The graph \(B=G_1\) adds \(p=c_0\) and \(x_0,x_1,x_2\) to \(A\).
The centre is adjacent to every vertex of \(A\) and to no auxiliary.
The neighbours of \(x_j\) in \(A\) are \(P_j\), and its auxiliary
adjacencies are specified by \(R\). Thus
\[
\card{V(B)}=20,\qquad \card{E(B)}=97+16+42+2=157.
\]
For the finite certificates, \(a_i\) has index \(i\), \(p=16\),
and \(x_j=17+j\). 
\subsection{The eight- and twelve-vertex prefixes}
The first two members are induced subgraphs of the same seed. Put
\[
\begin{aligned}
S_1&=\{2,4,5,8,9,10,12,13\},\\
S_2&=\{2,3,4,5,6,7,8,9,10,11,12,13\},
\end{aligned}
\]
and define
\[
H_1=A[\{a_i:i\in S_1\}],\qquad
H_2=A[\{a_i:i\in S_2\}],\qquad H_3=A.
\]
The direct reduction below establishes the strict level of the first prefix.
The recursive certificates then establish the strict level of the second
prefix and of the sixteen-vertex seed.
The two prefixes have 8 and 12 vertices, respectively.
The inclusions \(S_1\subset S_2\subset V(A)\) give induced inclusions.
We now verify the first prefix directly. Use numerical labels \(a_i=i\). Restricting the adjacency
table of \(A\) gives
\[
\begin{array}{c|l}
v&N_{H_1}(v)\\\hline
2&8,9,10,12,13\\
4&5,8,10,12,13\\
5&4,9,10,12,13\\
8&2,4,9,10,12\\
9&2,5,8,13\\
10&2,4,5,8,13\\
12&2,4,5,8\\
13&2,4,5,9,10
\end{array}
\]
\begin{proposition}\label{4k4:prop:manual-prefix}
The graph \(H_1\) belongs to \(\D_1\setminus\D_0\).
\end{proposition}
\begin{proof}
First consider \(H_1[N_{H_1}(5)]\), on \(\{4,9,10,12,13\}\).
The only neighbour of \(9\) there is \(13\), so \(9\to13\)
is a fold. The remaining graph on \(\{4,10,12,13\}\) is a cone
with apex \(4\). Consequently \(5\) can be deleted from \(H_1\)
at level one.

In \(H_1-5\), the neighbourhood of \(4\) is
\(\{8,10,12,13\}\). Within this neighbourhood, \(12\to8\)
is a fold, and the remaining path \(8,10,13\) is a cone with
apex \(10\). Thus \(4\) can next be deleted at level one.
The graph \(H_1-\{5,4\}\) is a cone with apex \(2\), proving
\(H_1\in\D_1\). The inner folds justify only the stated outer
deletion; they do not delete vertices from the outer graph.

To exclude level zero, Table~\ref{4k4:tab:manual-prefix-negative}
lists, for every \(v\), each candidate apex \(t\in N_{H_1}(v)\)
and a vertex \(w\in N_{H_1}(v)\setminus\{t\}\) not adjacent to it.
Every candidate occurs, so no neighbourhood has an apex.
There is no first fold, and \(H_1\notin\D_0\).
\end{proof}

\begin{table}[htbp]
\centering\small
\caption{Complete exclusion of an apex in every neighbourhood of \(H_1\).}
\label{4k4:tab:manual-prefix-negative}
\begin{tabular}{r l}\toprule
\(v\)&Candidate apex \(t\mapsto w\), with \(tw\notin E(H_1)\)\\\midrule
2 & $8\mapsto 13,\ 9\mapsto 10,\ 10\mapsto 9,\ 12\mapsto 9,\ 13\mapsto 8$\\
4 & $5\mapsto 8,\ 8\mapsto 5,\ 10\mapsto 12,\ 12\mapsto 10,\ 13\mapsto 8$\\
5 & $4\mapsto 9,\ 9\mapsto 4,\ 10\mapsto 9,\ 12\mapsto 9,\ 13\mapsto 12$\\
8 & $2\mapsto 4,\ 4\mapsto 2,\ 9\mapsto 4,\ 10\mapsto 9,\ 12\mapsto 9$\\
9 & $2\mapsto 5,\ 5\mapsto 2,\ 8\mapsto 5,\ 13\mapsto 8$\\
10 & $2\mapsto 4,\ 4\mapsto 2,\ 5\mapsto 2,\ 8\mapsto 5,\ 13\mapsto 8$\\
12 & $2\mapsto 4,\ 4\mapsto 2,\ 5\mapsto 2,\ 8\mapsto 5$\\
13 & $2\mapsto 4,\ 4\mapsto 2,\ 5\mapsto 2,\ 9\mapsto 4,\ 10\mapsto 9$\\
\bottomrule
\end{tabular}
\end{table}
\subsection{Certifying the prefixes and the seed}\label{4k4:app:certificates}
For the larger graphs, the neighbourhood reductions are themselves recursive. We record them as finite proof states so that every required membership and non-membership can be checked.
The main certificate contains 1,114 proof nodes for \(A\) and 184
for \(B-p\). The prefix certificates use the ordered vertex lists
\(S_1,S_2\) of this section.
All neighbourhoods and domination relations are computed in the
current induced subgraph.

\subsubsection{The recursive proof rule}
A proof state \((S,j)\), with \(j\geq-1\), asserts whether the
graph induced by \(S\) belongs to \(\D_j\).
There are five allowed forms of node.
\begin{enumerate}
\item The empty state is false at every level.
\item A non-empty cone is true at every level, with its apex supplied.
\item A non-empty non-cone is false at level \(-1\).
\item A true state with \(j\geq0\) specifies \(v\in S\), a true child
\((S\cap N(v),j-1)\), and a true child \((S\setminus\{v\},j)\).
\item A non-empty non-cone state \((S,j)\), with \(j\geq0\)
(and hence \(\card S\geq2\)), is false if it specifies, for every
\(v\in S\), a false child of at least one of those two forms.
\end{enumerate}
The checker verifies the exact subsets, levels, truth values and
child references, requires every vertex in a negative node to be
covered, and rejects cycles. Every child has fewer vertices.
It reconstructs each induced graph from the stated adjacency data,
tests every proposed apex, and checks the absence of an apex in each
negative non-cone state. A deleted vertex must belong to the current set.
The correctness of a checked proof graph therefore follows by
induction on the number of vertices, using Definition~\ref{def:hierarchy}
and Lemma~\ref{lem:inclusions} for cone states. By
the first-deletion equivalence~\eqref{eq:first}, any reduction of a state with
\(\card S\geq2\) begins with a vertex \(v\) such that
\((S\cap N(v),j-1)\) is true, and the remaining reduction shows
that \((S\setminus\{v\},j)\) is true. The false children of a
form-5 node exclude every such \(v\). This gives a finite certificate
of non-membership as well as membership, without relying on a bounded
unsuccessful search.

The three main roots are bound to
\[
(V(A),3,\mathrm{true}),\quad (V(A),2,\mathrm{false}),\quad
(V(B)\setminus\{p\},3,\mathrm{true}).
\]
The four prefix roots assert membership at levels \(1,2\) and
non-membership at levels \(0,1\), respectively.
The checker also verifies the prefix adjacencies against the
specified induced subgraphs of \(A\).
The seed and first extension are reconstructed from
Tables~\ref{4k4:tab:Aadj} and~\ref{4k4:tab:profiles}, the type graph
in~\eqref{4k4:eq:R}, and the six adjacency rules of
Section~\ref{sec:construction}.

The complete finite proof graphs are the files
\path{data/certificates/FINITE_CERTIFICATE.json.gz} and
\path{data/certificates/PREFIX_CERTIFICATES.json.gz}.
The seed and profiles are in \path{data/certificates/PARAMETERS.json};
the ordered prefix sets and their adjacencies are in
\path{data/certificates/PREFIX.json}. The checking rules are implemented in
\path{dismantling/certificates.py}. The verifier checks the supplied proof
states using integer and set operations; it does not search for reductions.

\subsubsection{The reductions and their obstructions}
For \(H_1\), delete \(a_5,a_4\) in that order; the remaining graph
is a cone with apex \(a_2\). Their current neighbourhoods belong
to \(\D_0\), as proved by the referenced child certificates.
For \(H_2\), delete \(a_6,a_5,a_4\); the current neighbourhoods
belong to \(\D_1\), and the remainder is a cone with apex \(a_2\).
For \(A=H_3\), delete \(a_0,a_1\); the current neighbourhoods belong
to \(\D_2\), and the remainder is a cone with apex \(a_{15}\).
For \(B-p\), the outer order is
\[
a_3,a_1,a_2,a_0,a_4,a_6,a_5,a_8,a_9,a_{10},a_{11}.
\]
Each current neighbourhood belongs to \(\D_2\), and the remainder
is a cone with apex \(a_7\).
These orders alone do not certify their neighbourhoods;
the complete child proof graphs do so by the preceding rules.

The negative root for \(A\notin\D_2\) has one obstruction for each
of its 16 vertices: every \(A[N_A(a_i)]\) is certified outside
\(\D_1\). The prefix negative roots similarly provide a complete
obstruction for every possible first deletion at their claimed
lower levels. A zero value of \(q\) is never used as a sufficient test of membership.

\subsection{The nineteen neighbourhood obstructions}
Table~\ref{4k4:tab:Bq} supplies all coefficients through degree seven
of every non-central neighbourhood in \(B\).
All higher degrees were computed and verified to vanish.
They are checked using the coefficient recurrence~\eqref{eq:clique-coefficients}.
The supplied checker uses this deletion recurrence to verify every
coefficient and alternating sum in the distributed certificate.
For example, the row for \(a_0\) gives
\(1-16+95-263+344-182+20=-1\).
All 19 sums are non-zero.
\begin{table}[H]
\centering\small
\begin{tabular}{r rrrrrrrr r}\toprule
\(v\)&\(c_0\)&\(c_1\)&\(c_2\)&\(c_3\)&\(c_4\)&\(c_5\)&\(c_6\)&\(c_7\)&\(q\)\\\midrule
\(a_{0}\) & 1 & 16 & 95 & 263 & 344 & 182 & 20 & 0 & $-1$\\
\(a_{1}\) & 1 & 16 & 95 & 262 & 340 & 179 & 20 & 0 & $-1$\\
\(a_{2}\) & 1 & 16 & 96 & 271 & 362 & 189 & 10 & 0 & $-7$\\
\(a_{3}\) & 1 & 16 & 95 & 261 & 335 & 172 & 16 & 0 & $-2$\\
\(a_{4}\) & 1 & 16 & 96 & 275 & 391 & 257 & 62 & 3 & $-1$\\
\(a_{5}\) & 1 & 16 & 97 & 281 & 402 & 262 & 60 & 3 & $-2$\\
\(a_{6}\) & 1 & 16 & 95 & 263 & 349 & 200 & 34 & 2 & $-2$\\
\(a_{7}\) & 1 & 16 & 97 & 279 & 388 & 232 & 39 & 1 & $-3$\\
\(a_{8}\) & 1 & 15 & 86 & 235 & 308 & 164 & 15 & 0 & $-4$\\
\(a_{9}\) & 1 & 16 & 95 & 263 & 344 & 183 & 19 & 0 & $-3$\\
\(a_{10}\) & 1 & 16 & 95 & 264 & 354 & 207 & 37 & 1 & $-1$\\
\(a_{11}\) & 1 & 14 & 75 & 194 & 249 & 143 & 27 & 2 & $-1$\\
\(a_{12}\) & 1 & 16 & 96 & 271 & 367 & 211 & 34 & 2 & $-2$\\
\(a_{13}\) & 1 & 15 & 86 & 237 & 323 & 199 & 41 & 1 & $-1$\\
\(a_{14}\) & 1 & 16 & 97 & 281 & 402 & 262 & 60 & 3 & $-2$\\
\(a_{15}\) & 1 & 16 & 96 & 272 & 375 & 230 & 48 & 3 & $-1$\\
\(x_0\) & 1 & 15 & 86 & 236 & 315 & 179 & 26 & 0 & $-2$\\
\(x_1\) & 1 & 15 & 85 & 228 & 296 & 167 & 27 & 0 & $-1$\\
\(x_2\) & 1 & 16 & 97 & 278 & 380 & 213 & 27 & 1 & $-3$\\
\bottomrule\end{tabular}
\caption{Clique coefficients of \(B[N_B(v)]\), for every non-central
vertex. All coefficients of degree greater than seven are zero.
The last column is their alternating sum.}
\label{4k4:tab:Bq}
\end{table}



\runcheck{4k4}
The output reports seven checked roots and nineteen clique rows. The root
checks certify $H_1\in\D_1\setminus\D_0$, $H_2\in\D_2\setminus\D_1$,
$A\in\D_3\setminus\D_2$, and $B-p\in\D_3$. Every coefficient in the
printed table is compared with both clique enumeration and the deletion
recurrence, including the vanishing of higher coefficients.

\subsection{The family at every level}
The certificates establish the hypotheses needed for the extension from the sixteen-vertex seed. Theorem~\ref{thm:template} then treats any number of layers.
\begin{theorem}\label{4k4:thm:family}
With the seed and profiles in this section, put $H_3=A$ and
$H_k=G_{k-3}$ for $k\geq3$, and use the two induced prefixes for $k=1,2$.
Then
\[
 H_k\in\D_k\setminus\D_{k-1},\qquad |V(H_k)|=4k+4\qquad(k\geq1).
\]
Every inclusion $H_k\subset H_{k+1}$ is induced.
\end{theorem}
\begin{proof}
The seven roots give the two prefixes and establish
$A\in\D_3\setminus\D_2$ and $B-p\in\D_3$. The nineteen non-zero sums
in Table~\ref{4k4:tab:Bq} exclude every non-central vertex as a first deletion.
Thus Theorem~\ref{thm:template} applies with $s=3$ and $d=3$.
For $k\geq3$ it gives strict level $k$ and
$16+4(k-3)=4k+4$ vertices. The prefix certificates prove the cases
$k=1,2$. The nested prefix sets and the layer rules give the induced inclusions.
\end{proof}
For $k\geq3$, the layered graphs also have
\[
 |E(G_m)|=97+16m+42m+4\binom m2+2m^2+3m(m-1)
          =97+53m+7m^2.
\]
The terms count seed edges, seed--centre edges, seed--auxiliary edges,
edges between centres or within one type, edges between the two adjacent
pairs of types, and centre--auxiliary edges, respectively.

\section{The family of order \texorpdfstring{$6k+6$}{6k+6}}\label{sec:family6}
The same layer construction also applies to a twelve-vertex seed at strict
level one. This example uses five auxiliary types, and its positive
certificates can be written out entirely as folds inside successive
neighbourhoods. It illustrates how the three finite hypotheses are checked
when no recursive certificate file is needed.
We first state the finite assertions and their consequences; the reductions
and clique tables used in the proof are displayed in
Section~\ref{6k6:app:certificates}.

\subsection{The seed and the resulting family}\label{6k6:sec:seed}
Here five auxiliary types give six new vertices per layer. Let
\(V(A)=\{a_0,\ldots,a_{11}\}\), with the adjacency list in
Table~\ref{6k6:tab:Aadj}. The list is symmetric, has no loops, and gives
\(\card{E(A)}=47\). The profiles are given in Table~\ref{6k6:tab:profiles};
their cardinalities sum to \(48\). The type graph is
\begin{equation}\label{6k6:eq:R}
 E(R)=\{02,03,04,12,14,23,34\},
\end{equation}
where \(jt\) abbreviates the unordered pair \(\{j,t\}\).
There are no other edges of \(R\).
\begin{table}[H]
\centering\small\setlength{\tabcolsep}{4pt}
\caption{The complete adjacency list of \(A\). An integer \(i\) denotes \(a_i\).}\label{6k6:tab:Aadj}
\begin{tabular}{c l}
\toprule
$v$ & $N_A(v)$\\
\midrule

0 & \(\{1,3,4,5,6,7,9,10\}\)\\
1 & \(\{0,2,3,4,5,6,7,8,10\}\)\\
2 & \(\{1,4,5,6,8,9,11\}\)\\
3 & \(\{0,1,4,6,7,8,9,11\}\)\\
4 & \(\{0,1,2,3,5,8,10,11\}\)\\
5 & \(\{0,1,2,4,6,9,10,11\}\)\\
6 & \(\{0,1,2,3,5,8,9\}\)\\
7 & \(\{0,1,3,9,10,11\}\)\\
8 & \(\{1,2,3,4,6,9,10,11\}\)\\
9 & \(\{0,2,3,5,6,7,8,10,11\}\)\\
10 & \(\{0,1,4,5,7,8,9,11\}\)\\
11 & \(\{2,3,4,5,7,8,9,10\}\)\\
\bottomrule
\end{tabular}
\end{table}

\begin{table}[H]
\centering\small\setlength{\tabcolsep}{4pt}
\caption{The five profiles. Each displayed integer \(i\) denotes \(a_i\).}\label{6k6:tab:profiles}
\begin{tabular}{c l}\toprule
$j$ & Indices of the vertices in $P_j$\\\midrule
0 & \(1,2,3,4,5,6,7,8,9,11\)\\
1 & \(1,3,4,6,7,8,10,11\)\\
2 & \(0,1,3,4,5,6,7,9,10,11\)\\
3 & \(1,2,3,5,6,7,8,9,10,11\)\\
4 & \(0,1,2,3,5,7,8,9,10,11\)\\
\bottomrule\end{tabular}
\end{table}

The graph \(B\) consists of \(A\), the centre \(p\), and the five
auxiliaries \(x_0,\ldots,x_4\). Its centre is adjacent to every vertex
of \(A\) and to no auxiliary; the neighbours of \(x_j\) in \(A\)
are precisely \(P_j\), and its auxiliary adjacencies are those of \(R\).
Thus
\[
 \card{V(B)}=18,\qquad
 \card{E(B)}=47+12+48+7=114.
\]
In the finite certificates, \(a_i\) is abbreviated to \(i\), and
\(p=12\), \(x_j=13+j\). In particular, vertex \(15\) denotes \(x_2\).
\begin{proposition}\label{6k6:prop:seed}
These seed data satisfy
\[
 A\in\D_1\setminus\D_0,\qquad B-p\in\D_1,\qquad
 q(B[N_B(v)])=
 \begin{cases}
 2,&v=a_7\text{ or }v=x_2,\\
 1,&v\ne p,a_7,x_2
 \end{cases}
 \quad(v\ne p).
\]
\end{proposition}
\begin{proof}
Table~\ref{6k6:tab:Ano0} gives, for every \(v\in V(A)\), the degree
\(d_v=\card{N_A(v)}\) and the maximum degree in \(A[N_A(v)]\).
In each row that maximum is strictly smaller than \(d_v-1\).
No open neighbourhood graph is a cone; Lemma~\ref{lem:first}
therefore gives \(A\notin\D_0\).

Tables~\ref{6k6:tab:Apos} and~\ref{6k6:tab:Bpos} specify successive outer
deletions in \(A\) and \(B-p\), respectively. Each row lists the
exact neighbourhood set in the current outer graph, the folds
within that neighbourhood in the stated order, and an apex of
the resulting neighbourhood. All fourteen required closed
neighbourhood inclusions are displayed in Table~\ref{6k6:tab:inclusiones}.
They are obtained by restricting the adjacency data of
Tables~\ref{6k6:tab:Aadj} and~\ref{6k6:tab:profiles} and~\eqref{6k6:eq:R} to the
current set. After the outer deletions, the remaining vertex sets
are cones with the apices specified in Section~\ref{6k6:app:certificates}.
Each inner reduction is consequently a level-\(0\) reduction, and
each outer deletion is permitted at level \(1\).
The final cones also lie in \(\D_1\), by Lemma~\ref{lem:inclusions}.
This proves \(A\in\D_1\) and \(B-p\in\D_1\).

Table~\ref{6k6:tab:Bq} gives all clique coefficients of the seventeen
non-central neighbourhood graphs. These coefficients follow by
the coefficient recurrence in Lemma~\ref{lem:qrec}, with
\(c_0(H)=1\) and \(c_t(\varnothing)=0\) for \(t\geq1\).
Their alternating sums are \(2\) in rows \(7\) and \(15\), and
\(1\) in every other row. This proves the last assertion.
\end{proof}

For this seed, the graphs of Section~\ref{sec:construction} satisfy
\begin{equation}\label{6k6:eq:concrete-order}
 \card{V(G_m)}=12+6m,\qquad
 \card{E(G_m)}=47+52m+15m^2 .
\end{equation}
Indeed, the seven disjoint classes of edges give
\[
\begin{aligned}
 \card{E(G_m)}
 &=47+12m+48m+\binom m2+5\binom m2+7m^2+5m(m-1)\\
 &=47+52m+15m^2.
\end{aligned}
\]
The terms count, in order, edges within \(A\), between \(A\) and
centres, between \(A\) and auxiliaries, between centres, within
each auxiliary type, between distinct types, and between centres
and auxiliaries.

\begin{theorem}\label{6k6:thm:family}
For every $k\geq1$, put $J_k=G_{k-1}$. Then $J_k\in\D_k\setminus\D_{k-1}$ and $|V(J_k)|=6k+6$.
\end{theorem}
\begin{proof}
Proposition~\ref{6k6:prop:seed} supplies the three hypotheses of
Theorem~\ref{thm:template}, with \(s=1\) and \(d=5\). Consequently
\[
 G_m\in\D_{m+1}\setminus\D_m\qquad(m\geq0).
\]
Taking \(m=k-1\) gives \(J_k\in\D_k\setminus\D_{k-1}\) and
\(\card{V(J_k)}=12+6(k-1)=6k+6\).
\end{proof}

\subsection{The finite reductions and obstructions}\label{6k6:app:certificates}
The certificates below use only the adjacency list, the profiles,
and the type edges specified in Section~\ref{6k6:sec:seed}.

\subsubsection{Excluding a first fold in \texorpdfstring{\(A\)}{A}}
For \(v\in V(A)\), put \(d_v=\card{N_A(v)}\), and write
\(\Delta(Q)\) for the maximum degree in a non-empty graph \(Q\).
The entries of Table~\ref{6k6:tab:Ano0} satisfy
\(\Delta(A[N_A(v)])<d_v-1\), so none of these neighbourhoods
has an apex.
\begin{table}[H]
\centering\small\setlength{\tabcolsep}{4pt}
\caption{Every open neighbourhood graph of \(A\) has no universal vertex.}\label{6k6:tab:Ano0}
\begin{tabular}{c r r r}
\toprule
$v$ & $d_v=\card{N_A(v)}$ & $d_v-1$ & $\Delta(A[N_A(v)])$\\
\midrule

0 & 8 & 7 & 6\\
1 & 9 & 8 & 6\\
2 & 7 & 6 & 5\\
3 & 8 & 7 & 5\\
4 & 8 & 7 & 6\\
5 & 8 & 7 & 5\\
6 & 7 & 6 & 5\\
7 & 6 & 5 & 4\\
8 & 8 & 7 & 5\\
9 & 9 & 8 & 6\\
10 & 8 & 7 & 5\\
11 & 8 & 7 & 6\\
\bottomrule
\end{tabular}
\end{table}


\subsubsection{The two level-1 reductions}
In each row of Tables~\ref{6k6:tab:Apos} and~\ref{6k6:tab:Bpos}, \(v\) is
deleted from the outer graph after all preceding outer deletions.
The set \(S\) is its neighbourhood set at that moment.
The arrows are folds inside the graph induced by \(S\).
After each arrow its source vertex is removed before the next
inclusion is evaluated. The listed apex is universal after
those inner folds. A row with no folds already has a conical
neighbourhood graph.

For Table~\ref{6k6:tab:Apos}, the final outer vertex set is
\(\{5,6,8,9,10,11\}\), with apex \(9\).
For Table~\ref{6k6:tab:Bpos}, it is
\(\{3,7,8,9,10,11,13,14,15,16,17\}\), with apex \(11\).
The centre \(p\) has been removed before Table~\ref{6k6:tab:Bpos}
begins.
\begin{table}[H]
\centering\small\setlength{\tabcolsep}{4pt}
\caption{A reduction certificate for \(A\in\D_1\).}\label{6k6:tab:Apos}
\begin{tabular}{c p{0.45\textwidth} p{0.22\textwidth} c}
\toprule
$v$ & Current neighbourhood set $S$ & Folds in order & Apex\\
\midrule

7 & \(\{0,1,3,9,10,11\}\) & \(1\to 0\) & 9\\
0 & \(\{1,3,4,5,6,9,10\}\) & \(4\to 1\), \(3\to 6\) & 5\\
1 & \(\{2,3,4,5,6,8,10\}\) & \(3\to 8\), \(6\to 2\) & 4\\
2 & \(\{4,5,6,8,9,11\}\) & \(4\to 11\) & 9\\
3 & \(\{4,6,8,9,11\}\) & none & 8\\
4 & \(\{5,8,10,11\}\) & none & 10\\
\bottomrule
\end{tabular}
\end{table}

\Needspace{16\baselineskip}
\begin{table}[H]
\centering\small\setlength{\tabcolsep}{4pt}
\caption{A reduction certificate for \(B-p\in\D_1\).}\label{6k6:tab:Bpos}
\begin{tabular}{c p{0.45\textwidth} p{0.22\textwidth} c}
\toprule
$v$ & Current neighbourhood set $S$ & Folds in order & Apex\\
\midrule

0 & \(\{1,3,4,5,6,7,9,10,15,17\}\) & \(4\to 1\), \(6\to 15\), \(3\to 7\) & 10\\
2 & \(\{1,4,5,6,8,9,11,13,16,17\}\) & none & 13\\
4 & \(\{1,3,5,8,10,11,13,14,15\}\) & \(5\to 15\), \(10\to 14\) & 3\\
1 & \(\{3,5,6,7,8,10,13,14,15,16,17\}\) & \(5\to 16\), \(6\to 3\), \(8\to 17\) & 7\\
5 & \(\{6,9,10,11,13,15,16,17\}\) & none & 9\\
6 & \(\{3,8,9,13,14,15,16\}\) & none & 3\\
\bottomrule
\end{tabular}
\end{table}

For example, the first row of Table~\ref{6k6:tab:Apos} starts with
\(S=\{0,1,3,9,10,11\}\). Within this neighbourhood,
\[
 N_{A[S]}[1]=\{0,1,3,10\}\subseteq
 N_{A[S]}[0]=\{0,1,3,9,10\}.
\]
The fold \(1\to0\) leaves a cone with apex \(9\), so the outer
vertex \(7\) can be deleted at level \(1\).
The next row is evaluated in \(A-7\).
Table~\ref{6k6:tab:inclusiones} provides the corresponding sets for
every fold in both certificates. Its third column is contained
in its fourth in each row; both columns are the closed
neighbourhoods in the current inner graph.
\begin{table}[H]
\centering\small\setlength{\tabcolsep}{4pt}
\caption{Closed neighbourhood inclusions for all fourteen folds. Here \(T\) is the neighbourhood graph after the preceding folds in the same row of the corresponding certificate.}\label{6k6:tab:inclusiones}
\begin{tabular}{c c p{0.31\textwidth} p{0.36\textwidth}}
\toprule
Graph; outer $v$ & Fold $u\to t$ & $N_T[u]$ & $N_T[t]$\\
\midrule

$A$; 7 & \(1\to 0\) & \(\{0,1,3,10\}\) & \(\{0,1,3,9,10\}\)\\
$A$; 0 & \(4\to 1\) & \(\{1,3,4,5,10\}\) & \(\{1,3,4,5,6,10\}\)\\
$A$; 0 & \(3\to 6\) & \(\{1,3,6,9\}\) & \(\{1,3,5,6,9\}\)\\
$A$; 1 & \(3\to 8\) & \(\{3,4,6,8\}\) & \(\{2,3,4,6,8,10\}\)\\
$A$; 1 & \(6\to 2\) & \(\{2,5,6,8\}\) & \(\{2,4,5,6,8\}\)\\
$A$; 2 & \(4\to 11\) & \(\{4,5,8,11\}\) & \(\{4,5,8,9,11\}\)\\
$B-p$; 0 & \(4\to 1\) & \(\{1,3,4,5,10,15\}\) & \(\{1,3,4,5,6,7,10,15,17\}\)\\
$B-p$; 0 & \(6\to 15\) & \(\{1,3,5,6,9,15\}\) & \(\{1,3,5,6,7,9,10,15\}\)\\
$B-p$; 0 & \(3\to 7\) & \(\{1,3,7,9,15,17\}\) & \(\{1,3,7,9,10,15,17\}\)\\
$B-p$; 4 & \(5\to 15\) & \(\{1,5,10,11,13,15\}\) & \(\{1,3,5,10,11,13,14,15\}\)\\
$B-p$; 4 & \(10\to 14\) & \(\{1,8,10,11,14,15\}\) & \(\{1,3,8,10,11,14,15\}\)\\
$B-p$; 1 & \(5\to 16\) & \(\{5,6,10,13,15,16,17\}\) & \(\{3,5,6,7,8,10,13,15,16,17\}\)\\
$B-p$; 1 & \(6\to 3\) & \(\{3,6,8,13,14,15,16\}\) & \(\{3,6,7,8,13,14,15,16,17\}\)\\
$B-p$; 1 & \(8\to 17\) & \(\{3,8,10,13,14,16,17\}\) & \(\{3,7,8,10,13,14,16,17\}\)\\
\bottomrule
\end{tabular}
\end{table}


\subsubsection{Clique coefficients}
Table~\ref{6k6:tab:Bq} supplies the coefficients \(c_0,\ldots,c_6\)
of every non-central neighbourhood graph in \(B\).
There are no cliques of greater cardinality in these graphs.
The coefficient recurrence~\eqref{eq:clique-coefficients}
computes every row from the adjacency data.
For instance, the alternating sums in rows \(0\) and \(7\) are
\[
 1-11+41-63+35-2=1,\qquad
 1-12+49-86+63-15+2=2.
\]
All seventeen sums are non-zero.
\begin{table}[H]
\centering\small\setlength{\tabcolsep}{4pt}
\caption{Clique counts for the seventeen non-central open neighbourhood graphs of \(B\). Coefficients \(\kappa_t\) with \(t>6\) are zero.}\label{6k6:tab:Bq}
\begin{tabular}{c r r r r r r r r}
\toprule
$v$ & $\kappa_0$ & $\kappa_1$ & $\kappa_2$ & $\kappa_3$ & $\kappa_4$ & $\kappa_5$ & $\kappa_6$ & $q$\\
\midrule

0 & 1 & 11 & 41 & 63 & 35 & 2 & 0 & 1\\
1 & 1 & 15 & 72 & 133 & 88 & 12 & 0 & 1\\
2 & 1 & 11 & 43 & 75 & 58 & 17 & 2 & 1\\
3 & 1 & 14 & 65 & 122 & 91 & 23 & 3 & 1\\
4 & 1 & 12 & 46 & 71 & 39 & 2 & 0 & 1\\
5 & 1 & 13 & 57 & 103 & 74 & 16 & 1 & 1\\
6 & 1 & 12 & 46 & 75 & 51 & 10 & 0 & 1\\
7 & 1 & 12 & 49 & 86 & 63 & 15 & 2 & 2\\
8 & 1 & 13 & 57 & 103 & 75 & 18 & 2 & 1\\
9 & 1 & 14 & 66 & 131 & 111 & 37 & 5 & 1\\
10 & 1 & 13 & 56 & 97 & 61 & 7 & 0 & 1\\
11 & 1 & 14 & 65 & 125 & 101 & 32 & 5 & 1\\
13 & 1 & 13 & 59 & 119 & 110 & 42 & 5 & 1\\
14 & 1 & 10 & 33 & 45 & 22 & 0 & 0 & 1\\
15 & 1 & 13 & 57 & 103 & 74 & 15 & 1 & 2\\
16 & 1 & 13 & 60 & 125 & 120 & 47 & 5 & 1\\
17 & 1 & 13 & 56 & 104 & 85 & 28 & 4 & 1\\
\bottomrule
\end{tabular}
\end{table}


\runcheck{6k6}
The output confirms twelve outer deletions, fourteen inner folds, twelve
rows excluding a first fold in the seed, and seventeen clique rows.
Each displayed neighbourhood and both sides of every fold inclusion are
recomputed from the adjacency lists. The calculation also verifies the
final cones in the two reductions.

\section{Completing kernels with two auxiliary types}\label{sec:kernels}
The family $F_k$ adds only two auxiliaries per layer. For this case we can
understand the neighbourhood obstructions by removing certain edges from the
seed. We first introduce the homology used to compare the resulting graphs,
then derive the identities and apply them to the eleven-vertex seed.

\subsection{Homology for the comparison}
A non-zero reduced homology group excludes a graph from every dismantling
class. We give the chain conventions and the elementary facts used in this
section and in Section~\ref{sec:homology}; in particular, empty
neighbourhoods must remain visible in degree $-1$.
The clique complex \(\Cl(H)\) contains the cliques of \(H\), including
the empty face. We use \defi{augmented reduced homology over \(\F\)}. For a complex \(K\), the chain group
\(\Ct_t(K)\) has the faces of cardinality \(t+1\) as a basis when
\(t\geq0\), and
\(\Ct_{-1}(K)=\F\), generated by the empty face. Lower chain groups
are zero. The boundary of a face is the sum of its codimension-one
faces; in particular the boundary of a vertex is the empty face.
Each term occurs twice in \(\partial^2\), so this is a chain complex.

We write \(\Ht_t(K)=\ker\partial_t/\operatorname{im}\partial_{t+1}\)
and call \(K\) \defi{acyclic} if every \(\Ht_t(K)\) vanishes.
In particular,
\begin{equation}\label{hom:eq:empty}
 \Ht_{-1}(\{\varnothing\})=\F,\qquad
 \Ht_t(\{\varnothing\})=0\quad(t\ne-1).
\end{equation}
A complex with a vertex has zero homology in degree \(-1\).
Its \defi{reduced Betti number} in degree \(t\) is
\(\widetilde\beta_t(K)=\dim_{\F}\Ht_t(K)\).
Acyclicity is used as a homological condition, not as a synonym for
contractibility or for membership in any \(\D_j\).
When a homological term is applied to a graph, it refers to its clique
complex; in particular, \emph{acyclic} here does not mean that the graph is a forest.

A \defi{simplicial cone} with apex \(w\) is a complex
\(w*J\) whose faces are exactly the faces \(\sigma\) of \(J\)
and the sets \(\sigma\cup\{w\}\), where \(w\notin V(J)\).
In particular, if a graph \(H\) is a cone with apex \(w\), then
\(\Cl(H)=w*\Cl(H-w)\).

\begin{lemma}\label{hom:lem:cone}
A simplicial cone is acyclic.
\end{lemma}
\begin{proof}
For an apex \(w\), let \(h\) adjoin \(w\) to a face not containing it
and send a face containing \(w\) to zero. Direct expansion gives
\(\partial h+h\partial=\mathrm{id}\), including on the empty face.
Every cycle \(z\) is therefore the boundary of \(h(z)\).
\end{proof}

We will compare complexes by expressing them as unions of acyclic pieces.

\begin{lemma}\label{hom:lem:union}
If \(K=U\cup V\) and \(U,V\) are acyclic, then for every integer \(t\),
\[
 \Ht_t(K)\cong\Ht_{t-1}(U\cap V).
\]
The intersection may be \(\{\varnothing\}\).
\end{lemma}
\begin{proof}
For \(t\geq0\), write a cycle \(z=u+v\), with chains in \(U,V\).
Then \(\partial u=\partial v=b\) is a cycle in \(U\cap V\).
Changing the decomposition adds the same chain of the intersection
to \(u\) and \(v\), so it changes \(b\) by a boundary. If \(z\) is
changed by a boundary, decompose a filling into chains of \(U,V\);
the resulting class of \(b\) is unchanged. This defines
\([z]\mapsto[b]\).

Every cycle \(b\) of the intersection has fillings \(u\) in \(U\)
and \(v\) in \(V\), since both complexes are acyclic; then \(u+v\)
proves surjectivity. If \(b=\partial a\) in the intersection, the
chains \(u+a,v+a\) are cycles in \(U,V\) and hence boundaries.
Their sum is \(z\), proving injectivity.
For \(t=-1\), the left side is zero since the acyclic complexes
\(U,V\) have vertices; the right side is in degree \(-2\).
In lower degrees both sides vanish.
\end{proof}

The folds used in the graph reductions preserve these homology groups.

\begin{lemma}\label{hom:lem:fold}
A fold \(v\to w\) induces isomorphisms
\(\Ht_t(\Cl(H-v))\cong\Ht_t(\Cl(H))\) for every integer \(t\).
\end{lemma}
\begin{proof}
Pair every face \(\sigma\) containing \(v\) but not \(w\) with
\(\tau=\sigma\cup\{w\}\). The domination condition guarantees that
\(\tau\) is a face. Remove these pairs in decreasing order of
\(|\sigma|\). When a pair is removed, \(\tau\) is the only proper
coface of \(\sigma\) still present: every larger coface was removed
in an earlier pair. The remaining faces form a subcomplex.

To see that removing one such pair preserves augmented homology,
let \(K'\) be the resulting subcomplex of the current complex \(K\)
and put \(r=\dim\sigma\). A cycle of degree \(r+1\) has coefficient
zero at \(\tau\), since its boundary has coefficient zero at
\(\sigma\). For a cycle of degree \(r\), adding its coefficient at
\(\sigma\) times \(\partial\tau\) removes \(\sigma\) without changing
its homology class. In other degrees no basis element is removed.
Thus every class has a representative in \(K'\).

Conversely, suppose a cycle \(z'\) in \(K'\) bounds a chain \(y\)
in \(K\). If \(y\) has degree \(r+1\), its coefficient at \(\tau\)
is zero, because the coefficient of \(\sigma\) in \(z'\) is zero.
If \(y\) has degree \(r\), add its coefficient at \(\sigma\) times
\(\partial\tau\); this removes \(\sigma\) and preserves its
boundary. The resulting filling lies in \(K'\). This proves
injectivity as well, including in degree \(-1\).
After all pairs are removed the complex is exactly \(\Cl(H-v)\).
\end{proof}

The recursive neighbourhood condition therefore gives a homological obstruction at every level.

\begin{lemma}\label{hom:lem:acyclic}
If \(H\in\D_j\), \(j\geq-1\), then \(\Cl(H)\) is acyclic.
\end{lemma}
\begin{proof}
Cones are covered by Lemma~\ref{hom:lem:cone}. Induct on \(j\) and,
at a fixed level, on the number of vertices. The one-vertex graph
is a cone. For a valid first deletion \(v\), the two induction
hypotheses give acyclicity of \(\Cl(H-v)\) and \(\Cl(L_H(v))\).
The exact decomposition
\[
 \Cl(H)=\Cl(H-v)\cup\bigl(v*\Cl(L_H(v))\bigr)
\]
has acyclic pieces and intersection \(\Cl(L_H(v))\).
Lemma~\ref{hom:lem:union} proves the assertion.
\end{proof}



\subsection{The extension and its kernel}
In a construction with two auxiliary types joined by an edge, the profiles
determine which seed edges contribute to the homological obstructions.
Let \(A\) be a non-empty graph, and choose \(S,T\subseteq V(A)\).
Form \(B\) by adding vertices \(p,x,y\). The centre \(p\) is adjacent
to all of \(A\) and to neither auxiliary vertex. The neighbours of \(x\) in \(A\) are exactly \(S\), those of \(y\) are exactly \(T\), and \(xy\) is an edge.
The induced graph on \(A\) remains unchanged.

\begin{lemma}\label{ker:lem:cover}
If every non-central neighbourhood of \(B\) has non-trivial
reduced homology, then \(S\cup T=V(A)\).
The same conclusion holds if \(q(L_B(v))\ne0\) for every
\(v\ne p\).
\end{lemma}
\begin{proof}
If \(a\notin S\cup T\), then its neighbours in \(B\) are \(N_A(a)\)
and \(p\). The centre is an apex of \(L_B(a)\).
This cone has zero reduced homology and \(q=0\), contradicting
either condition.
\end{proof}

Henceforth assume \(S\cup T=V(A)\), and put
\[
 S_0=S\setminus T,\qquad T_0=T\setminus S,\qquad Z=S\cap T.
\]
Define \(H\) by deleting from \(A\) all its edges between \(S_0\)
and \(T_0\), and no others. We call \(H\) the \defi{kernel of these profile data}. This name describes an operation on edges; it does
not impose a homological property on \(H\).

Set
\[
 K=\Cl(A),\quad X=\Cl(A[S]),\quad Y=\Cl(A[T]),\quad
 W=\Cl(A[Z]),\quad U=\Cl(H).
\]
Every clique of \(H\) is contained in \(S\) or in \(T\), since it
cannot meet both \(S_0\) and \(T_0\). Conversely every clique contained
in \(S\) or in \(T\) survives the edge deletions. Therefore
\begin{equation}\label{ker:eq:kernel}
 U=X\cup Y,\qquad W=X\cap Y.
\end{equation}
In particular this union is a clique complex, even if it is smaller
than \(K\).

\subsection{The two auxiliary stars}

To compare the seed, its extension and the remainder, we group the faces containing auxiliary vertices. For a vertex \(z\) outside a complex \(J\), \(z*J\) denotes its cone.
Define the subcomplex
\begin{equation}\label{ker:eq:V}
 V=(x*X)\cup\bigl(y*(Y\cup(x*W))\bigr).
\end{equation}
It consists of all faces using an auxiliary vertex together with
their subfaces. Its two displayed pieces are cones and their
intersection is \(x*W\), also a cone. This intersection contains
the vertex \(x\), even when \(W=\{\varnothing\}\).
By Lemma~\ref{hom:lem:union}, \(V\) is acyclic.
The faces of \(V\) on the original vertices are exactly \(X\cup Y\).
Consequently
\begin{equation}\label{ker:eq:decompositions}
 \Cl(B)=(p*K)\cup V,\qquad
 \Cl(B-p)=K\cup V,\qquad
 (p*K)\cap V=K\cap V=U.
\end{equation}
These are equalities of complexes.

\begin{proposition}\label{ker:prop:global}
For every integer \(t\),
\[
 \Ht_t(\Cl(B))\cong\Ht_{t-1}(U).
\]
If \(K\) is acyclic, then also
\[
 \Ht_t(\Cl(B-p))\cong\Ht_{t-1}(U).
\]
Under this extra hypothesis, \(\Cl(B-p)\) is acyclic if and only if
\(U\) is acyclic.
\end{proposition}
\begin{proof}
The two pieces in the first decomposition of
\eqref{ker:eq:decompositions} are acyclic and intersect in \(U\).
Apply Lemma~\ref{hom:lem:union}. For the second decomposition the same
argument applies when \(K\) is acyclic. The shift covers every
possible non-zero degree; hence the last equivalence follows.
\end{proof}

For a seed \(A\in\D_s\), membership of \(B-p\) in a dismantling
class therefore implies acyclicity of the kernel. The converse
statement about membership does not follow: the proposition
characterises acyclicity of the remainder, not its dismantling level.

\subsection{Original and auxiliary neighbourhoods}

The global decomposition also applies within each neighbourhood. For an original vertex, its intersection is the neighbourhood complex in the kernel.

\begin{theorem}\label{ker:thm:original-links}
For every \(a\in V(A)\) and every integer \(t\),
\[
 \Ht_t(\Cl(L_B(a)))
 \cong \Ht_{t-1}(\Cl(L_H(a))).
\]
\end{theorem}
\begin{proof}
Put \(K_a=\Cl(A[N_A(a)])\).
If \(a\in S_0\), its neighbourhood in \(B\) contains \(p,x\) but
not \(y\), and
\[
 \Cl(L_B(a))=(p*K_a)\cup(x*X_a),\qquad
 X_a=\Cl(A[N_A(a)\cap S]).
\]
Both pieces are cones and their intersection is \(X_a\).
The edge deletions defining \(H\) remove exactly the neighbours
of \(a\) in \(T_0\). Edges among its remaining neighbours in \(S\)
are unchanged. Thus \(X_a=\Cl(L_H(a))\), and
Lemma~\ref{hom:lem:union} proves the assertion. The case \(a\in T_0\)
is obtained by exchanging the two profiles and auxiliaries.

If \(a\in Z\), none of the edges incident to \(a\) is deleted.
The neighbourhood contains both auxiliaries. Restrict the
construction \eqref{ker:eq:V} to \(N_A(a)\), writing
\[
 X_a=\Cl(A[N_A(a)\cap S]),\quad
 Y_a=\Cl(A[N_A(a)\cap T]),\quad
 W_a=X_a\cap Y_a .
\]
The resulting complex \(V_a\) is acyclic by the same union-of-cones
argument. Now \(\Cl(L_B(a))=(p*K_a)\cup V_a\), with intersection
\(X_a\cup Y_a\). The proof of \eqref{ker:eq:kernel}, restricted to
\(N_A(a)\), identifies this intersection with \(\Cl(L_H(a))\).
Lemma~\ref{hom:lem:union} again applies.
Augmented chains cover all possible empty neighbourhoods.
\end{proof}

The auxiliary neighbourhoods have the exact descriptions
\begin{equation}\label{ker:eq:auxiliary-complexes}
 \Cl(L_B(x))=X\cup(y*W),\qquad
 \Cl(L_B(y))=Y\cup(x*W).
\end{equation}
They do not contain \(p\).

\begin{lemma}\label{ker:lem:intersection-map}
If \(U=X\cup Y\) is acyclic, the inclusions induce isomorphisms
\[
 \Ht_t(W)\ \longrightarrow\ \Ht_t(X)\oplus\Ht_t(Y)
\]
for all integers \(t\).
\end{lemma}
\begin{proof}
For surjectivity, let \(a,b\) be cycles in \(X,Y\).
Since \(a+b\) is a cycle in \(U\), it bounds a chain
\(w=w_X+w_Y\) there. Then
\(z=a+\partial w_X=b+\partial w_Y\) is a cycle in \(W\)
representing the prescribed pair of homology classes.
For injectivity, suppose a cycle \(z\) in \(W\) bounds chains
\(u\) in \(X\) and \(v\) in \(Y\). The cycle \(u+v\) bounds
\(w=w_X+w_Y\) in \(U\). The chain
\(u+\partial w_X=v+\partial w_Y\) lies in \(W\) and has
boundary \(z\). These arguments apply in degree \(-1\) as well,
using the augmentation, and lower degrees are zero.
\end{proof}

When the kernel is acyclic, the auxiliary neighbourhoods are controlled by the opposite profile subgraphs.

\begin{theorem}\label{ker:thm:auxiliary-links}
If \(U\) is acyclic, then for every integer \(t\),
\[
 \Ht_t(\Cl(L_B(x)))\cong\Ht_{t-1}(Y),\qquad
 \Ht_t(\Cl(L_B(y)))\cong\Ht_{t-1}(X).
\]
\end{theorem}
\begin{proof}
By Lemma~\ref{ker:lem:intersection-map}, the map
\(\Ht_r(W)\to\Ht_r(X)\) is surjective and its kernel is
isomorphic to \(\Ht_r(Y)\): under the lemma's isomorphism
it is the projection onto the first summand.

For \(t\geq0\), write a cycle of \(X\cup(y*W)\) as \(u+y*w\).
The equations for being a cycle are
\(\partial w=0\) and \(\partial u=w\).
Assign to it the class of \(w\), which lies in the kernel of
\(\Ht_{t-1}(W)\to\Ht_{t-1}(X)\).
The assignment respects boundaries: the boundary of \(a+y*b\)
changes \(w\) by \(\partial b\).
Every class in that kernel has a filling \(u\) in \(X\), so
the assignment is surjective.

If \(w=\partial d\) in \(W\), adding
\(\partial(y*d)=d+y*\partial d\) leaves the cycle \(u+d\) in \(X\).
Surjectivity of \(\Ht_t(W)\to\Ht_t(X)\) lets us replace this
cycle, modulo a boundary in \(X\), by a cycle in \(W\).
That cycle bounds in the cone \(y*W\).
The original cycle therefore bounds in \(X\cup(y*W)\),
proving injectivity. Its homology in degree \(-1\) is zero
because it contains \(y\); the claimed right side then has
degree \(-2\) and is also zero. Lower degrees vanish.
Use \eqref{ker:eq:auxiliary-complexes} and exchange \(X,Y\)
for the second identity.
\end{proof}

\subsection{Equivalent finite conditions and their Euler form}

Here the layer construction of Section~\ref{sec:construction} has
$d=2$, $R=K_2$ and profiles $S,T$. We may therefore write the auxiliaries
as $x_i,y_i$, and $|V(G_m)|=|V(A)|+3m$. The identities just proved separate
the obstructions on original vertices from those on the two auxiliary types.

\begin{theorem}\label{ker:thm:equivalence}
Assume \(S\cup T=V(A)\), \(s\geq0\),
\(A\in\D_s\setminus\D_{s-1}\), and \(B-p\in\D_{s+1}\).
Then the following are equivalent:
\begin{enumerate}
\item Every non-central neighbourhood \(L_B(v)\), \(v\ne p\),
has non-trivial reduced homology.
\item The clique complex of the kernel \(H\) is acyclic, every \(L_H(a)\),
\(a\in V(A)\), has non-trivial reduced homology, and both
\(A[S]\) and \(A[T]\) have non-trivial reduced homology.
\end{enumerate}
Either condition implies
\[
 G_m\in\D_{s+m}\setminus\D_{s+m-1}\quad(m\geq0).
\]
\end{theorem}
\begin{proof}
The seed and remainder hypotheses imply acyclicity of \(K\) and
\(\Cl(B-p)\). Proposition~\ref{ker:prop:global} then implies
acyclicity of \(U=\Cl(H)\).
Theorem~\ref{ker:thm:original-links} gives equivalence of the
conditions on original vertices. Theorem~\ref{ker:thm:auxiliary-links}
gives equivalence of those on the auxiliaries.
Under condition (1), the data satisfy precisely the homological
criterion in Theorem~\ref{hom:thm:criterion}, proved independently in
Section~\ref{sec:homology}. It applies to
the displayed layered construction and proves strictness at
every subsequent level.
\end{proof}

The remainder hypothesis is present on both sides of this
equivalence. Acyclicity of the kernel is redundant under that
hypothesis, but states explicitly the property needed when
starting with a kernel. The equivalence does not dispense with the
remainder hypothesis.
Lemma~\ref{ker:lem:cover} gives the covering condition on the profiles.

Extend \(q\) from graphs to finite complexes by putting
\(q(J)=\sum_{\sigma\in J}(-1)^{|\sigma|}\), including the empty
face. This agrees with the existing notation \(q(H)=q(\Cl(H))\).
Inclusion--exclusion and the zero value on cones give identities
without homological assumptions.

\begin{proposition}\label{ker:prop:euler}
For the profile data above,
\begin{align}
 q(B)&=-q(H),&
 q(B-p)&=q(A)-q(H),\label{ker:eq:global-q}\\
 q(L_B(a))&=-q(L_H(a)) &&(a\in V(A)),\label{ker:eq:fixed-q}\\
 q(L_B(x))&=q(A[S])-q(A[Z]),&
 q(L_B(y))&=q(A[T])-q(A[Z]).\label{ker:eq:aux-q}
\end{align}
If \(q(H)=0\), the last line becomes
\[
 q(L_B(x))=-q(A[T]),\qquad
 q(L_B(y))=-q(A[S]).
\]
\end{proposition}
\begin{proof}
In \eqref{ker:eq:V}, both pieces and their intersection are cones,
so \(q(V)=0\). Apply inclusion--exclusion to
\eqref{ker:eq:decompositions} to obtain \eqref{ker:eq:global-q}.
The decompositions used to prove
Theorem~\ref{ker:thm:original-links} likewise have pieces of
zero \(q\) and intersection \(\Cl(L_H(a))\), yielding
\eqref{ker:eq:fixed-q}.
Apply inclusion--exclusion to \eqref{ker:eq:auxiliary-complexes}
to obtain \eqref{ker:eq:aux-q}.
Finally \eqref{ker:eq:kernel} gives
\(q(H)=q(A[S])+q(A[T])-q(A[Z])\).
\end{proof}

In particular, under \(A\in\D_s\setminus\D_{s-1}\) and
\(B-p\in\D_s\), both have \(q=0\), so \(q(H)=0\).
The non-zero Euler conditions on all non-central neighbourhoods
of \(B\) are then equivalent to
\[
 q(L_H(a))\ne0\quad\hbox{for all }a\in V(A),\qquad
 q(A[S])\ne0,\quad q(A[T])\ne0.
\]
This is the kernel formulation of Theorem~\ref{thm:template}. The identities remain valid if the remainder
is only assumed to lie in \(\D_{s+1}\), as in its separate
relaxation. The equality \(q=0\) is not used as a
test for acyclicity or dismantlability.

One further necessary condition can be read from
Lemma~\ref{ker:lem:intersection-map}. If \(S,T,Z\) are non-empty
and \(H\) is acyclic, the numbers of connected components satisfy
\[
 c(A[Z])=c(A[S])+c(A[T])-1.
\]
Indeed, take dimensions in degree zero and use
\(\dim\Ht_0=c-1\).
This numerical equality alone does not imply acyclicity; even
agreement of dimensions in all degrees would not establish that
the inclusion maps in the lemma are isomorphisms.

\subsection{Completing one edge and reducing the remainder}

The case of a single added edge lets us give an explicit reduction of the remainder. Suppose \(H\) is a graph, \(ab\) is a non-edge, and
\[
 A=H+ab,\qquad S=V(H)\setminus\{b\},\qquad
 T=V(H)\setminus\{a\}.
\]
Here \(S_0=\{a\}\) and \(T_0=\{b\}\), so deleting the edges between
the exclusive profile parts recovers exactly \(H\).

\begin{lemma}\label{ker:lem:delete-vertex}
If \(\Cl(H)\) is acyclic, inclusion induces
\[
 \Ht_t(\Cl(L_H(v)))\cong\Ht_t(\Cl(H-v))
\]
for every vertex \(v\) and every integer \(t\).
\end{lemma}
\begin{proof}
Write \(P=\Cl(H-v)\), \(Q=v*\Cl(L_H(v))\) and
\(I=P\cap Q=\Cl(L_H(v))\). The union is acyclic by hypothesis,
and \(Q\) is a cone. The same chain argument as in
Lemma~\ref{ker:lem:intersection-map}, now with the acyclic union
\(P\cup Q=\Cl(H)\), shows that
\(\Ht_t(I)\to\Ht_t(P)\oplus\Ht_t(Q)\) is an isomorphism.
The second summand is zero.
\end{proof}

Thus, if \(H\) is acyclic and all its neighbourhoods have
non-trivial reduced homology, then both \(A[S]=H-b\) and
\(A[T]=H-a\) have non-trivial reduced homology.
The remaining obligations for
Theorem~\ref{ker:thm:equivalence} are the strict level of the completed
seed and the required level of \(B-p\).

There is a useful sufficient condition for the latter.
For disjoint graphs \(C,J\), write \(C\star J\) for their
\defi{complete join}.

\begin{lemma}\label{ker:lem:join}
If \(C\in\D_j\), \(j\geq-1\), then
\(C\star J\in\D_j\) for every graph \(J\).
\end{lemma}
\begin{proof}
For \(j=-1\), an apex of \(C\) remains universal.
Induct on \(j\). Follow a level-\(j\) reduction of \(C\).
The neighbourhood of a deleted vertex in the join is the join
of its current neighbourhood in \(C\) with \(J\), and belongs
to \(\D_{j-1}\) by induction. When only one vertex of \(C\)
remains, the join is a cone and reduces at level zero, hence at
level \(j\). This proves the assertion for \(j\geq0\).
\end{proof}

\begin{proposition}\label{ker:prop:remainder-folds}
For the one-edge completion, put
\[
 C=H[N_H(a)\cap N_H(b)].
\]
The remainder \(B-p\) folds onto \(C\star C_4\), where the cycle
has vertices \(a,b,y,x\) in this order.
Consequently, if \(C\in\D_j\) for some \(j\geq0\), then
\(B-p\in\D_j\).
\end{proposition}
\begin{proof}
Every \(z\in V(H)\setminus\{a,b\}\) that is not adjacent to \(a\) in \(H\)
is dominated by \(y\) in \(B-p\). Indeed \(y\) is adjacent to
every vertex except \(a\), and \(a\) is not in the closed
neighbourhood of \(z\). Likewise, a vertex not adjacent to \(b\)
is dominated by \(x\).
Delete all vertices outside \(C\cup\{a,b,x,y\}\), retaining
\(a,b,x,y\) and the common neighbours. Each chosen domination
inclusion remains true after restricting to the vertices still
present.

In the remaining graph, the edges on \(\{a,b,x,y\}\) are exactly
\(ab,by,yx,xa\), and every vertex of \(C\) is adjacent to these
four vertices. Its induced graph is unchanged from \(H\).
The remainder is therefore \(C\star C_4\); when \(C\) is empty,
it is just the four-cycle.
If \(C\in\D_j\), Lemma~\ref{ker:lem:join} gives
\(C\star C_4\in\D_j\). Prepending the displayed folds is valid
at level \(j\) by nesting of the classes.
\end{proof}

If \(C\) is a cone, the proposition applies with \(j=0\).
The argument gives a sufficient certificate of the remainder.
It does not assert that a failure to dismantle \(C\) at a chosen
level excludes membership of \(B-p\) at that level.
In particular, an unsuccessful sufficient test must not replace
the original remainder condition in the equivalence theorem.

\subsection{The completion underlying the eleven-vertex seed}

Define the completed seed \(A\) on \(0,\ldots,10\) by
Table~\ref{ker:tab:seed}. It has thirty-three edges.
The pair \(\{1,2\}\) is an edge; define
\(V(H)=V(A)\) and \(E(H)=E(A)\setminus\{\{1,2\}\}\).
Then \(H\) has thirty-two edges and eleven vertices.
For the profiles \(S=V(A)\setminus\{2\}\) and
\(T=V(A)\setminus\{1\}\), the exclusive parts are
\(\{1\}\) and \(\{2\}\), respectively. Hence this \(H\) is
exactly the kernel of the profile data, and adding the single edge
\(\{1,2\}\) to \(H\) gives \(A\).

\begin{table}[htbp]
\centering
\caption{The complete adjacency data of the completed seed \(A\).}
\label{ker:tab:seed}\label{tab:seed-adjacency}
\begin{tabular}{cl}
\toprule
\(v\)&\(N_A(v)\)\\\midrule
$0$ & $\{3,4,6,7,8\}$\\
$1$ & $\{2,4,5,6,7,9,10\}$\\
$2$ & $\{1,5,6,7,8,9,10\}$\\
$3$ & $\{0,6,7,8,9,10\}$\\
$4$ & $\{0,1,7,8,9,10\}$\\
$5$ & $\{1,2,6,8,9\}$\\
$6$ & $\{0,1,2,3,5,10\}$\\
$7$ & $\{0,1,2,3,4,9\}$\\
$8$ & $\{0,2,3,4,5,10\}$\\
$9$ & $\{1,2,3,4,5,7\}$\\
$10$ & $\{1,2,3,4,6,8\}$\\
\bottomrule
\end{tabular}

\end{table}

This is the completed seed $F_1$ used in the article. The reductions of
Appendix~\ref{sec:article} establish its strict level and dismantle $B-p$.
The homology calculation gives
\[
 \Ht_*(\Cl(H))=0,\qquad
 \dim_{\F}\Ht_1(\Cl(L_H(v)))=
 \begin{cases}2,&v=10,\\1,&v\ne10,\end{cases}
\]
with all other reduced homology groups of these neighbourhoods zero.
Moreover, $N_H(1)\cap N_H(2)=\{5,6,7,9,10\}$ induces the path
$10,6,5,9,7$. Proposition~\ref{ker:prop:remainder-folds} therefore gives
$B-p\in\D_0$. Explicitly, the folds $0,3,8\to y$ followed by $4\to x$
leave the join of this path with the four-cycle on $1,2,y,x$.
The displayed neighbourhood groups and Lemma~\ref{ker:lem:delete-vertex}
give condition~(2) of Theorem~\ref{ker:thm:equivalence}.
With the strict level of $A$ proved in Appendix~\ref{sec:article}, that
theorem provides a homological proof of the strict levels of $F_k$.
\runcheck{kernels}
The report reconstructs $H$ by removing $\{1,2\}$, restores the seed by adding
that edge, computes the displayed homology groups over $\F$, and checks the
common-neighbour graph. It also compares the homology of each neighbourhood
in $H$ with the degree-shifted group in $B$.

\section{Homology through successive layers}\label{sec:homology}
The kernel calculations compare neighbourhoods in the first extension.
We now follow those groups through every layer of the general construction,
with any number of auxiliary types. As in Section~\ref{sec:construction},
write $\mathcal{T}=\{0,\ldots,d-1\}$ for the type set. The coefficients and augmented chain
conventions are those of Section~\ref{sec:kernels}. The degree shift will
supply the criterion used in Theorem~\ref{ker:thm:equivalence}.

\subsection{A degree shift in every non-central neighbourhood}

For complexes on disjoint vertex sets, \(K*J\) consists of unions
of one face from each factor. If \(S^0\) consists of two vertices
without an edge, \(K*S^0\) is the union of two cones on \(K\).
Lemma~\ref{hom:lem:union} gives
\begin{equation}\label{hom:eq:suspension}
 \Ht_t(K*S^0)\cong\Ht_{t-1}(K).
\end{equation}
This includes \(K=\{\varnothing\}\), for which the join is \(S^0\).
The clique complex of the complete join of two graphs is the join
of their clique complexes, since a clique is a union of a clique
in each graph.

\begin{theorem}\label{hom:thm:shift}
For \(m\geq1\) and every integer \(t\),
\begin{align}
 \Ht_t\bigl(\Cl(L_{G_m}(a))\bigr)
 &\cong\Ht_{t-m+1}\bigl(\Cl(L_B(a))\bigr),
 &&a\in V(A),\label{hom:eq:fixed}\\
 \Ht_t\bigl(\Cl(L_{G_m}(x_{j,i}))\bigr)
 &\cong\Ht_{t-m+1}\bigl(\Cl(L_B(x_j))\bigr),
 &&j\in \mathcal{T},\quad 0\leq i<m.\label{hom:eq:auxiliary}
\end{align}
\end{theorem}
\begin{proof}
First fix \(a\in V(A)\) and put \(\Lambda_m=L_{G_m}(a)\).
Its vertices are \(N_A(a)\), all centres, and all copies of the
types \(j\) for which \(a\in P_j\). For \(m\geq2\), delete a centre
\(c_i\) from \(\Lambda_m\). The folds \(x_{j,h}\to x_{j,i}\), \(h\ne i\),
are still valid: restrict the neighbourhood inclusions in the proof of
Lemma~\ref{lem:remainder} to \(N_A(a)\) and the types present.
After these folds, the remaining graph is a cone with a remaining
centre as apex.
If there are no such auxiliary types, it is already a cone.
Lemmas~\ref{hom:lem:cone} and \ref{hom:lem:fold} imply that
\(\Cl(\Lambda_m-c_i)\) is acyclic.

The neighbourhood of \(c_i\) inside \(\Lambda_m\) is isomorphic to
\(\Lambda_{m-1}\), using the same layer reindexing restricted to the
vertices under consideration. Hence the decomposition
\[
 \Cl(\Lambda_m)=\Cl(\Lambda_m-c_i)
       \cup\bigl(c_i*\Cl(L_{\Lambda_m}(c_i))\bigr)
\]
has acyclic pieces and intersection isomorphic to \(\Cl(\Lambda_{m-1})\).
Lemma~\ref{hom:lem:union} shifts homology by one degree. Iteration
proves \eqref{hom:eq:fixed}; at \(m=1\) it is the identity.
The set \(N_A(a)\) may be empty, but \(\Lambda_1\) still contains \(p\).

Now fix an auxiliary vertex \(x_{j,i}\). Its neighbourhood has
the vertices of \(P_j\), all copies of types in \(N_R(j)\),
the vertices \(x_{j,h}\) for \(h\ne i\), and the centres \(c_h\)
for \(h\ne i\). For each \(u\in N_R(j)\), fold
\(x_{u,h}\) onto \(x_{u,i}\) for \(h\ne i\). The source and
target have the same adjacencies to \(P_j\) and auxiliary types.
The target sees all centres in this neighbourhood, whereas
the source may omit \(c_h\). These folds therefore preserve
homology, and their targets remain throughout.

What remains consists of a copy of \(L_B(x_j)\) on
\[
 P_j\ \sqcup\ \{x_{u,i}:u\in N_R(j)\}
\]
and the \(m-1\) pairs \(\{c_h,x_{j,h}\}\), \(h\ne i\).
Each pair has no internal edge. Distinct pairs are completely
joined, and every pair is completely joined to the displayed copy
of \(L_B(x_j)\). The remaining clique complex is consequently
\[
 \Cl(L_B(x_j))*(S^0)^{*(m-1)}.
\]
Apply Lemma~\ref{hom:lem:fold} and then \eqref{hom:eq:suspension}
\(m-1\) times. If \(L_B(x_j)\) is empty, both \(P_j\) and
\(N_R(j)\) are empty. There are then no folds, the first factor is
\(\{\varnothing\}\), and the same suspension identity applies.
For \(m=1\) there are no pairs, so this case too is the identity.
\end{proof}

\subsection{The homological criterion}

Non-zero groups in the first layer remain non-zero after every shift. Together with a reduction of the remainder, they give the following criterion for strictness.

\begin{theorem}\label{hom:thm:criterion}
Let \(s\geq0\). Suppose that
\begin{enumerate}
\item \(A\in\D_s\setminus\D_{s-1}\);
\item \(B-p\in\D_{s+1}\);
\item for every \(v\in V(B)\setminus\{p\}\), some integer \(t\)
satisfies \(\Ht_t(\Cl(L_B(v)))\ne0\).
\end{enumerate}
Then for every \(m\geq0\),
\[
 G_m\in\D_{s+m}\setminus\D_{s+m-1},
 \qquad |V(G_m)|=|V(A)|+(d+1)m.
\]
\end{theorem}
\begin{proof}
Theorem~\ref{hom:thm:shift} propagates non-zero reduced homology to
every non-central neighbourhood in every \(G_m\), \(m\geq1\).
By Lemma~\ref{hom:lem:acyclic}, none of these neighbourhoods belongs
to any class \(\D_j\).

The case \(m=0\) is the first hypothesis. For \(m=1\), the centre
has neighbourhood \(A\in\D_s\) and its deletion leaves
\(B-p\in\D_{s+1}\). Equation~\eqref{eq:first} yields
\(B\in\D_{s+1}\). A first deletion at level \(s\) is impossible:
the centre's neighbourhood is not in \(\D_{s-1}\), and neither is
any non-central neighbourhood. Since \(B\) has more than one
vertex, \eqref{eq:first} gives \(B\notin\D_s\).
This includes \(s=0\), when \(\D_{s-1}\) is the class of cones.

For \(m\geq2\), suppose \(G_{m-1}\) is strictly
\((s+m-1)\)-dismantlable. Each centre has neighbourhood isomorphic
to \(G_{m-1}\); deleting it leaves a graph in
\(\D_0\subseteq\D_{s+m}\) by Lemmas~\ref{lem:centre} and~\ref{lem:remainder}.
Thus \(G_m\in\D_{s+m}\). A first deletion at level \(s+m-1\)
would require a neighbourhood in \(\D_{s+m-2}\).
Centres fail this condition by the induction hypothesis, and
non-central vertices fail it by homology. This proves strictness.
The order follows from \eqref{eq:general-order}.
\end{proof}

\subsection{Comparison with Euler characteristic}

The obstruction \(q\) of Section~\ref{sec:construction} is the alternating
sum of reduced Betti numbers, with a contribution from degree \(-1\).
Rank--nullity in the augmented chain complex gives
\begin{equation}\label{hom:eq:euler}
 q(H)=\sum_{r\geq0}(-1)^r
       \dim_{\F}\Ht_{r-1}(\Cl(H)).
\end{equation}
Indeed, if \(f_t=\dim\Ct_t\), \(b_t=\dim\Ht_t\), and
\(\rho_t=\operatorname{rank}\partial_t\), then
\(f_t=b_t+\rho_t+\rho_{t+1}\). In the finite alternating sum
the ranks cancel, including the term of degree \(-1\).
In particular, \(q(H)\ne0\) implies non-trivial reduced homology.

The converse need not hold. The graph \(C_4\sqcup K_1\) has five
vertices, four edges and no triangles, so \(q=1-5+4=0\).
Its two components give \(\dim\Ht_0=1\). The boundary matrix of
the four cycle edges has rank three, and there are no
two-dimensional faces, giving \(\dim\Ht_1=1\).
These contributions cancel in \eqref{hom:eq:euler}.
This example separates the two obstructions on a graph; it does
not by itself supply initial data satisfying all three hypotheses
of Theorem~\ref{hom:thm:criterion}.

\Needspace{10\baselineskip}
The Euler criterion in Theorem~\ref{thm:template} uses
\(q(L_B(v))\ne0\) for every non-central \(v\), together with
\(A\in\D_s\setminus\D_{s-1}\) and the stronger remainder condition
\(B-p\in\D_s\). Equation~\eqref{hom:eq:euler} and
Lemma~\ref{lem:inclusions} imply the hypotheses of Theorem~\ref{hom:thm:criterion}.
Allowing the remainder in \(\D_{s+1}\) is a separate change from
replacing Euler characteristic by homology: the first-deletion
argument at \(m=1\) also permits this weaker remainder in the
Euler proof. Both criteria retain a distinct requirement on the
remainder; non-trivial homology of the other neighbourhoods alone
does not ensure the positive reduction.


\runcheck{homology}
The calculation checks the empty graph, a cone, $C_4\sqcup K_1$, and an
auxiliary with an empty first-layer neighbourhood. In the last example,
after $m$ layers its only non-zero reduced Betti number is $\widetilde\beta_{m-2}=1$,
for the tested values $m=1,2,3$. The full command
\texttt{python reproduce.py --extended} additionally checks all labelled
seeds on one, two or three vertices, all profiles for one or two types,
and all type graphs, at $m=0,1,2,3$. Homology is computed for the subcollection
with seed order at most two. These are finite checks of the identities;
Theorem~\ref{hom:thm:shift} proves them for arbitrary seeds and all layers.

\appendix
\section{The article's eleven-vertex seed and its extension}\label{sec:article}\label{sec:example}\label{sec:seed3k8}
The seed $F_1=A$ is given in Table~\ref{ker:tab:seed}. Here we record the
finite reductions and Euler obstructions that establish the strict levels
of the $3k+8$ family.
Its profiles and type graph are
\begin{equation}\label{eq:seed-profiles}
 P_0=V(A)\setminus\{2\},\qquad P_1=V(A)\setminus\{1\},\qquad R=K_2.
\end{equation}
The first extension $F_2=B$ adds $p,x,y$, with $p$ adjacent to every seed
vertex and to neither auxiliary. The vertex $x$ is adjacent to $P_0$, the
vertex $y$ to $P_1$, and $xy$ is an edge. Thus $A$ has eleven vertices and
thirty-three edges, while $B$ has fourteen vertices and sixty-five edges.
Figure~\ref{fig:extension} shows the added edges in this first extension.

\begin{figure}[htbp]
\centering\begin{tikzpicture}[x=1cm,y=1cm]
\coordinate (w) at (0,0);
\coordinate (p) at (0,2);
\coordinate (x) at (-3,.6);
\coordinate (y) at (3,.6);
\coordinate (a) at (-3,-1.6);
\coordinate (b) at (3,-1.6);
\draw[jfef new edge] (p)--(w);
\draw[jfef new edge] (x)--(w);
\draw[jfef new edge] (y)--(w);
\draw[jfef new edge] (x)--(a);
\draw[jfef new edge] (y)--(b);
\draw[jfef new edge] (x) .. controls (-1.7,-1.7) and (1.7,-1.7) .. (y);
\draw[jfef new edge] (p) .. controls (-5,2.2) and (-5,-1.2) .. (a);
\draw[jfef new edge] (p) .. controls (5,2.2) and (5,-1.2) .. (b);
\node[ellipse,draw=black!70,fill=white,minimum width=3cm,minimum height=1.3cm,jfef label font] at (w) {$W$};
\node[jfef centre] at (p) {$p$};
\node[jfef auxiliary] at (x) {$x$};
\node[jfef auxiliary] at (y) {$y$};
\node[jfef vertex] at (a) {$1$};
\node[jfef vertex] at (b) {$2$};
\end{tikzpicture}

\caption{The thirty-two added edges. The ellipse $W$ represents
$V(A)\setminus\{1,2\}$, the set $Z$ of Section~\ref{sec:kernels}; an edge to $W$ means adjacency to all nine vertices.
Edges within $A$, including $\{1,2\}$, are specified by Table~\ref{ker:tab:seed}.}\label{fig:extension}
\end{figure}

\Needspace{6\baselineskip}
\subsection{The seed is strictly \texorpdfstring{$1$}{1}-dismantlable}
We will use folds in two different graphs. To show that \(A\) belongs
to level one, we fold each neighbourhood graph needed in a reduction.
To exclude level zero, we examine the folds available in \(A\) itself.

\begin{lemma}\label{lem:seed-in-D1}
The graph \(A\) belongs to \(\D_1\).
\end{lemma}
\begin{proof}
We delete the vertices of \(A\) in the order
\[
 5,2,1,4,0,6,7,8,9,3,
\]
leaving \(10\). Each row of Table~\ref{tab:seed-reduction} justifies
one deletion. The second column lists the neighbours of the vertex to
be deleted, after the preceding deletions from \(A\). The third column
gives the folds in the graph induced by those neighbours, and the last
column gives the surviving vertex.

To check a listed fold \(u\to t\), let \(H\) be the neighbourhood
graph for that row after its preceding folds. The adjacencies in
Table~\ref{tab:seed-adjacency} give \(N_H[u]\subseteq N_H[t]\), so the
fold is valid by Lemma~\ref{lem:dom}. In the final two rows the
neighbourhood graph is already a single vertex. Thus each row gives a
reduction of its neighbourhood graph to one vertex, which puts that
graph in \(\D_0\) by Definition~\ref{def:hierarchy}. The ten deletions
therefore dismantle \(A\) at level one.
\end{proof}

\begin{table}[htbp]
\centering\small\setlength{\tabcolsep}{4pt}
\caption{A complete level-one reduction of \(A\).}\label{tab:seed-reduction}
\begin{tabular}{cclc}
\toprule
Deleted vertex&Current neighbourhood&Folds within that neighbourhood&Survivor\\\midrule
$5$&$\{1,2,6,8,9\}$&$1\to2,\ 6\to2,\ 8\to2,\ 2\to9$&$9$\\
$2$&$\{1,6,7,8,9,10\}$&$6\to1,\ 7\to1,\ 8\to10,\ 9\to1,\ 1\to10$&$10$\\
$1$&$\{4,6,7,9,10\}$&$6\to10,\ 7\to4,\ 9\to4,\ 4\to10$&$10$\\
$4$&$\{0,7,8,9,10\}$&$9\to7,\ 7\to0,\ 0\to8,\ 8\to10$&$10$\\
$0$&$\{3,6,7,8\}$&$6\to3,\ 7\to3,\ 3\to8$&$8$\\
$6$&$\{3,10\}$&$3\to10$&$10$\\
$7$&$\{3,9\}$&$3\to9$&$9$\\
$8$&$\{3,10\}$&$3\to10$&$10$\\
$9$&$\{3\}$&\text{none}&$3$\\
$3$&$\{10\}$&\text{none}&$10$\\
\bottomrule
\end{tabular}

\end{table}
\FloatBarrier

\Needspace{7\baselineskip}
\begin{lemma}\label{lem:seed-not-D0}
The graph \(A\) does not belong to \(\D_0\).
\end{lemma}
\begin{proof}
To determine whether a level-zero reduction can start, we list the
folds available in \(A\). By Lemma~\ref{lem:dom}, a fold of \(u\) onto
\(t\) in the current graph \(H\) is permitted exactly when \(u\ne t\)
and \(N_H[u]\subseteq N_H[t]\). Comparing the closed neighbourhoods from
Table~\ref{tab:seed-adjacency} gives the complete lists below, first
in \(A\) and then after each forced deletion:
\[
\begin{array}{c|c}
 H&\text{all permitted folds}\\\hline
 A&5\to2\\
 A-5&9\to7\\
 A-\{5,9\}&\text{none}.
\end{array}
\]
Any level-zero reduction would have to delete \(5\) and then \(9\).
The remaining graph on \(\{0,1,2,3,4,6,7,8,10\}\) admits no fold,
so the reduction stops before reaching \(K_1\).
\end{proof}

\subsection{A dismantlable remainder}
The neighbourhood graph of \(p\) is \(A\), so \(p\) can be deleted at
level two. Theorem~\ref{thm:template} uses the hypothesis that the
remaining graph \(B-p\) can be dismantled at level one. In fact,
ordinary folds already suffice.

\begin{lemma}\label{lem:seed-rest}
The graph \(B-p\) belongs to \(\D_0\).
\end{lemma}
\begin{proof}
After removing \(p\), the remaining vertices are \(0,\ldots,10,x,y\).
Perform the following twelve folds in order:
\begin{equation}\label{eq:seed-rest-folds}
\begin{aligned}
 &0\to x,\quad3\to x,\quad4\to x,\quad7\to9,\quad8\to y,\quad9\to5,\\
 &1\to6,\quad2\to6,\quad5\to6,\quad6\to10,\quad10\to x,\quad x\to y.
\end{aligned}
\end{equation}
At each step, the adjacency table and the specified adjacencies
of \(x,y\) give the closed-neighbourhood inclusion for the indicated
fold in the graph left by the preceding folds. Only \(y\) survives.
Thus \(B-p\in\D_0\subseteq\D_1\), giving the required hypothesis of
Theorem~\ref{thm:template} for \(s=1\).
\end{proof}

\Needspace{5\baselineskip}
\subsection{Euler obstructions in the first extension}
A first deletion at level one would require a neighbourhood graph
in \(\D_0\). The centre fails this test because its neighbourhood
graph is \(A\). For each other vertex, we use Lemma~\ref{lem:qobstruction}:
a non-zero \(q\) excludes its neighbourhood graph from \(\D_0\).

\begin{lemma}\label{lem:seed-euler}
Every non-central vertex \(v\) of \(B\) satisfies
\(q(B[N_B(v)])\ne0\).
\end{lemma}
\begin{proof}
For this clique table only, write \(H_v=B[N_B(v)]\) for each of the thirteen vertices \(v\ne p\). This local notation follows the article table and is unrelated to the family \(H_k\) of Section~\ref{sec:family4}.
To evaluate \(q(H_v)\), we count the cliques of \(H_v\) by size.
In Table~\ref{tab:seed-cliques}, \(c_r(H_v)\) counts the cliques of
size \(r\). The empty clique gives \(c_0(H_v)=1\), and none of the
neighbourhood graphs has a clique of size greater than four.

To obtain a row directly from the adjacency data, enumerate the
subsets of \(N_B(v)\) and count those in which every pair of vertices
is adjacent. These counts are the coefficients of \(C_{H_v}\).
Alternatively, compute the same polynomial using the recurrence in
Lemma~\ref{lem:qrec}, starting from \(C_{\varnothing}(z)=1\).
Each step uses graphs with fewer vertices, so the recursion terminates.
In either case, evaluation at \(z=-1\) gives the last column of
Table~\ref{tab:seed-cliques}, where every entry is non-zero.
\end{proof}

\begin{table}[htbp]
\centering
\caption{All clique counts of \(H_v=B[N_B(v)]\) for \(v\ne p\).}\label{tab:seed-cliques}
\begin{tabular}{crrrrrr}
\toprule
\(v\)&\(\kappa_0(B_v)\)&\(\kappa_1(B_v)\)&\(\kappa_2(B_v)\)&\(\kappa_3(B_v)\)&\(\kappa_4(B_v)\)&\(q(B_v)\)\\\midrule
$0$&1&8&21&20&5&$-1$\\
$1$&1&9&25&24&6&$-1$\\
$2$&1&9&24&21&4&$-1$\\
$3$&1&9&25&24&6&$-1$\\
$4$&1&9&25&24&6&$-1$\\
$5$&1&8&20&16&2&$-1$\\
$6$&1&9&24&21&4&$-1$\\
$7$&1&9&25&25&7&$-1$\\
$8$&1&9&24&21&4&$-1$\\
$9$&1&9&25&24&6&$-1$\\
$10$&1&9&25&23&4&$-2$\\
$x$&1&11&35&37&11&$-1$\\
$y$&1&11&35&36&10&$-1$\\
\bottomrule
\end{tabular}

\end{table}

\subsection{The family at every strict level}
The checks on \(A\) and \(B\) now let us apply Theorem~\ref{thm:template}
to any number of layers, including zero layers.

\begin{theorem}\label{thm:family3k8}
For every integer \(k\geq1\), the graph \(F_k=G_{k-1}\)
with \(d=2\) and \(R=K_2\), specified by Table~\ref{tab:seed-adjacency},
\eqref{eq:seed-profiles}, and the adjacency rules of
Section~\ref{sec:construction} satisfies
\[
 F_k\in\D_k\setminus\D_{k-1},\qquad |V(F_k)|=3k+8.
\]
\end{theorem}
\begin{proof}
Lemmas~\ref{lem:seed-in-D1} and~\ref{lem:seed-not-D0} establish
\(A\in\D_1\setminus\D_0\). The second hypothesis follows from
\(B-p\in\D_0\subseteq\D_1\) by Lemma~\ref{lem:seed-rest}, and the
last is exactly Lemma~\ref{lem:seed-euler}.

Apply Theorem~\ref{thm:template} with \(s=1\) and \(d=2\).
It gives strict level \(m+1\) and order \(11+3m\) after \(m\) layers,
for every \(m\geq0\). Taking \(m=k-1\) gives level \(k\) and order
\(11+3(k-1)=3k+8\). The first two graphs are \(F_1=A\) and \(F_2=B\).
\end{proof}

\runcheck{article}
The report checks all ten outer deletions and twenty-three inner folds of the
seed, the forced folds excluding level zero, the twelve folds of $B-p$, and
all thirteen clique rows. It also checks the exact incidences drawn in the
figures: the eight schematic incidences in Figure~\ref{fig:extension} expand to thirty-two added edges. Figure~\ref{fig:complement} displays the seed by its twenty-two missing edges; Figure~\ref{fig:layers} shows all thirty edges on the nine vertices added in three layers.

\begin{figure}[htbp]
\centering\begin{tikzpicture}[x=1cm,y=1cm]
\path[use as bounding box] (-3.35,-2.9) rectangle (3.35,2.9);
\coordinate (v0) at (0.000000,2.350000);
\coordinate (v1) at (1.270506,1.976946);
\coordinate (v2) at (2.137635,0.976225);
\coordinate (v3) at (2.326080,-0.334440);
\coordinate (v4) at (1.776011,-1.538923);
\coordinate (v5) at (0.662072,-2.254808);
\coordinate (v6) at (-0.662072,-2.254808);
\coordinate (v7) at (-1.776011,-1.538923);
\coordinate (v8) at (-2.326080,-0.334440);
\coordinate (v9) at (-2.137635,0.976225);
\coordinate (v10) at (-1.270506,1.976946);
\draw[jfef complement edge] (v0) -- (v1);
\draw[jfef complement edge] (v0) -- (v2);
\draw[jfef complement edge] (v0) -- (v5);
\draw[jfef complement edge] (v0) -- (v9);
\draw[jfef complement edge] (v0) -- (v10);
\draw[jfef complement edge] (v1) -- (v3);
\draw[jfef complement edge] (v1) -- (v8);
\draw[jfef complement edge] (v2) -- (v3);
\draw[jfef complement edge] (v2) -- (v4);
\draw[jfef complement edge] (v3) -- (v4);
\draw[jfef complement edge] (v3) -- (v5);
\draw[jfef complement edge] (v4) -- (v5);
\draw[jfef complement edge] (v4) -- (v6);
\draw[jfef complement edge] (v5) -- (v7);
\draw[jfef complement edge] (v5) -- (v10);
\draw[jfef complement edge] (v6) -- (v7);
\draw[jfef complement edge] (v6) -- (v8);
\draw[jfef complement edge] (v6) -- (v9);
\draw[jfef complement edge] (v7) -- (v8);
\draw[jfef complement edge] (v7) -- (v10);
\draw[jfef complement edge] (v8) -- (v9);
\draw[jfef complement edge] (v9) -- (v10);
\node[jfef vertex] at (v0) {$0$};
\node[jfef vertex] at (v1) {$1$};
\node[jfef vertex] at (v2) {$2$};
\node[jfef vertex] at (v3) {$3$};
\node[jfef vertex] at (v4) {$4$};
\node[jfef vertex] at (v5) {$5$};
\node[jfef vertex] at (v6) {$6$};
\node[jfef vertex] at (v7) {$7$};
\node[jfef vertex] at (v8) {$8$};
\node[jfef vertex] at (v9) {$9$};
\node[jfef vertex] at (v10) {$10$};
\end{tikzpicture}

\caption{The complement of the eleven-vertex seed, with the twenty-two
non-edges listed by the adjacency data.}\label{fig:complement}
\end{figure}
\begin{figure}[htbp]
\centering\begin{tikzpicture}[x=1cm,y=1cm]
\path[use as bounding box] (-3.1,-2.7) rectangle (3.1,2.7);
\coordinate (v0) at (-2.100000,2.000000);
\coordinate (v1) at (0.000000,2.000000);
\coordinate (v2) at (2.100000,2.000000);
\coordinate (v3) at (-2.100000,-0.000000);
\coordinate (v4) at (-2.100000,-2.000000);
\coordinate (v5) at (0.000000,-0.000000);
\coordinate (v6) at (0.000000,-2.000000);
\coordinate (v7) at (2.100000,-0.000000);
\coordinate (v8) at (2.100000,-2.000000);
\draw[jfef edge] (v0) -- (v1);
\draw[jfef edge] (v0) .. controls (-0.700000,2.900000) and (0.700000,2.900000) .. (v2);
\draw[jfef edge] (v0) -- (v5);
\draw[jfef edge] (v0) -- (v6);
\draw[jfef edge] (v0) -- (v7);
\draw[jfef edge] (v0) .. controls (-0.079310,1.318391) and (1.320690,-0.014943) .. (v8);
\draw[jfef edge] (v1) -- (v2);
\draw[jfef edge] (v1) -- (v3);
\draw[jfef edge] (v1) -- (v4);
\draw[jfef edge] (v1) -- (v7);
\draw[jfef edge] (v1) -- (v8);
\draw[jfef edge] (v2) -- (v3);
\draw[jfef edge] (v2) .. controls (1.320690,0.014943) and (-0.079310,-1.318391) .. (v4);
\draw[jfef edge] (v2) -- (v5);
\draw[jfef edge] (v2) -- (v6);
\draw[jfef edge] (v3) -- (v4);
\draw[jfef edge] (v3) -- (v5);
\draw[jfef edge] (v3) -- (v6);
\draw[jfef edge] (v3) .. controls (-0.700000,0.900000) and (0.700000,0.900000) .. (v7);
\draw[jfef edge] (v3) -- (v8);
\draw[jfef edge] (v4) -- (v5);
\draw[jfef edge] (v4) -- (v6);
\draw[jfef edge] (v4) -- (v7);
\draw[jfef edge] (v4) .. controls (-0.700000,-1.100000) and (0.700000,-1.100000) .. (v8);
\draw[jfef edge] (v5) -- (v6);
\draw[jfef edge] (v5) -- (v7);
\draw[jfef edge] (v5) -- (v8);
\draw[jfef edge] (v6) -- (v7);
\draw[jfef edge] (v6) -- (v8);
\draw[jfef edge] (v7) -- (v8);
\node[jfef centre] at (v0) {$c_{0}$};
\node[jfef centre] at (v1) {$c_{1}$};
\node[jfef centre] at (v2) {$c_{2}$};
\node[jfef auxiliary] at (v3) {$x_{0,0}$};
\node[jfef auxiliary] at (v4) {$x_{1,0}$};
\node[jfef auxiliary] at (v5) {$x_{0,1}$};
\node[jfef auxiliary] at (v6) {$x_{1,1}$};
\node[jfef auxiliary] at (v7) {$x_{0,2}$};
\node[jfef auxiliary] at (v8) {$x_{1,2}$};
\end{tikzpicture}

\caption{The new vertices in three layers. Each column contains its centre
and two auxiliaries. The seed and its incident edges are omitted. A centre
is adjacent to every other centre and to the auxiliaries outside its own
column. All six auxiliaries are pairwise adjacent.}\label{fig:layers}
\end{figure}

\subsection{Deleting one vertex}\label{sec:consequences}
Can a single vertex deletion preserve the strict level of \(F_k\)?
The proof of strictness already provides two ways to examine the
remaining graph. After deleting a centre we have an explicit fold
reduction. For any other vertex, the Euler recurrence relates the value
of \(q\) on the remaining graph to the value on its neighbourhood graph.

\begin{proposition}\label{prop:single-deletion}
Let $k\geq2$ and let $F_k$ be the graph of
Theorem~\ref{thm:family3k8}. Deleting a centre gives a graph in $\D_0$.
Deleting any other vertex gives a graph outside every class $\D_j$,
$j\geq-1$. Consequently no single vertex deletion preserves its
strict level $k$.
\end{proposition}
\begin{proof}
Write \(F_k=G_m\) with \(m=k-1\geq1\). If we delete a centre
when \(m=1\), the remainder is \(B-p\), which folds to one vertex by
Lemma~\ref{lem:seed-rest}. For \(m\geq2\), Lemma~\ref{lem:remainder}
gives a fold reduction after deleting any centre. In both cases
the remainder belongs to \(\D_0\subseteq\D_{k-1}\) and thus cannot
have strict level \(k\).

Now let \(v\) be a vertex other than a centre. By
Lemma~\ref{lem:euler-links}, the non-zero values in
Table~\ref{tab:seed-cliques} remain non-zero as we add layers, so
\(q(F_k[N_{F_k}(v)])\ne0\). The graph \(F_k\) itself belongs to
\(\D_k\) by Theorem~\ref{thm:family3k8}, so
\(q(F_k)=0\) by Lemma~\ref{lem:qobstruction}. Using these two
values in the recurrence of Lemma~\ref{lem:qrec} gives
\[
 q(F_k-v)=q(F_k)+q(F_k[N_{F_k}(v)])
         =q(F_k[N_{F_k}(v)])\ne0.
\]
Lemma~\ref{lem:qobstruction} therefore excludes \(F_k-v\) from
every dismantling class.
\end{proof}

The conclusion concerns one deletion from \(F_k\), starting with
\(F_2=B\). Minimality among all proper induced subgraphs and least
order among all graphs of strict level \(k\) would require separate
arguments.


\runcheck{deletions}
The calculation checks every single-vertex deletion from $F_2,F_3,F_4$.
It verifies the fold reduction after a centre deletion and the non-zero
Euler value after any other deletion. Proposition~\ref{prop:single-deletion}
proves the dichotomy for every $k\geq2$.

\FloatBarrier
\bibliographystyle{siam}
\bibliography{refs}
\end{document}
